\documentclass[11pt,a4paper]{article}
\usepackage[T1]{fontenc}
\usepackage{lmodern,microtype,amsmath,amssymb,amsthm,mathtools,array,booktabs,xcolor}
\usepackage[margin=25mm,top=24mm,bottom=24mm,headheight=15pt]{geometry}
\usepackage{fancyhdr,titlesec}
\definecolor{ink}{HTML}{173B47}\definecolor{accent}{HTML}{176C70}
\usepackage[colorlinks=true,linkcolor=ink,citecolor=accent,urlcolor=accent]{hyperref}
\usepackage{bookmark}
\titleformat{\section}{\large\sffamily\bfseries\color{ink}}{\thesection}{.5em}{}
\titleformat{\subsection}{\normalsize\sffamily\bfseries\color{ink}}{\thesubsection}{.5em}{}
\newtheorem{theorem}{Theorem}[section]\newtheorem{lemma}[theorem]{Lemma}\newtheorem{proposition}[theorem]{Proposition}\newtheorem{corollary}[theorem]{Corollary}
\numberwithin{equation}{section}
\newcommand{\R}{\mathbb R}\newcommand{\supp}{\operatorname{supp}}\newcommand{\vol}{\operatorname{vol}}
\pagestyle{fancy}\fancyhf{}\lhead{\small\sffamily Gaussian smoothing and extreme rays}\rhead{\small\sffamily Positive definite functions}\cfoot{\small\thepage}
\renewcommand{\headrulewidth}{.3pt}\setlength{\emergencystretch}{2em}\raggedbottom
\hypersetup{pdftitle={Gaussian smoothing and extreme rays of positive definite functions},pdfsubject={Fourier zeros, scale mixtures, convolution roots and spherical means}}
\begin{document}
\begin{center}
{\LARGE\bfseries\color{ink} Gaussian smoothing and extreme rays\\[3pt] of positive definite functions\par}
\vspace{7pt}{\large Fourier zeros, the Riemann $\Xi$-function, and scale mixtures}\par
\vspace{8pt}{\small Research manuscript prepared for Keio \enspace$\cdot$\enspace 20 November 2026}
\end{center}
\begin{abstract}
We answer the product and convolution questions posed by Jaming, Matolcsi and R\'ev\'esz in 2008. Both closure assertions fail in every dimension. On the line, the convolution of the tent with the standard Gaussian generates an extreme ray. Multiplying the result by the same Gaussian destroys extremality. For every continuous compactly supported radial nonnegative positive definite core in any dimension, its Gaussian smoothing is extreme precisely when its entire radial Fourier transform has only real zeros. If that transform has $4k$ nonreal zeros counted with multiplicity, the smallest radial face containing the smoothing has dimension $2k+1$. Every such dimension occurs. For the classical Riemann function $\Xi$ we show that $\Xi^2$ is extreme exactly when the Riemann hypothesis holds and that each off-line zero, counted with multiplicity, contributes one dimension when their number is finite. This criterion leaves the Riemann hypothesis undecided. We give criteria for the scale mixtures in the authors' Question~2 and establish Gaussian uniqueness for two classes of eventually log-concave tails. For compactly supported radial covariances we obtain sharp universal counts of real and complex half-ball convolution roots. An application to Efimov's spherical extremal problem gives a finite equation at every radius in odd dimensions and certified spectral bounds in all dimensions $d\ge2$.
\end{abstract}
\section*{Questions and scope}
Write $T(x)=(1-|x|)_+$, $\gamma(x)=e^{-\pi x^2}$, and call a nonzero member of a cone \emph{extreme} if every decomposition into two cone members has proportional summands. The cone in the first three rows below consists of continuous radial functions that are both nonnegative and positive definite. The questions and their one-dimensional examples come from Section~6 of \cite{JMR2009}.
\begin{center}\small\begin{tabular}{@{}p{.20\linewidth}p{.71\linewidth}@{}}\toprule
Original question & Result in this manuscript \\ \midrule
1(a), 1(b) & Both operations admit counterexamples in every dimension (Section~\ref{sec:jmr}).\\
2(a), 2(b), 2(c) & Ball autocorrelation and quartic Hermite scale-mixture tests, followed by a fixed-dimension criterion valid for signed functions (Sections~\ref{sec:ballmix}--\ref{sec:signed}).\\
3 & $T*\gamma$ is extreme and $\gamma(T*\gamma)$ admits two nonproportional summands (Section~\ref{sec:jmr}).\\
Borisov's question & The compact-core Gaussian class and the tails in Theorem~\ref{thm:tail} are covered. The general one-dimensional strict-positivity conjecture remains open here.\\ \bottomrule
\end{tabular}\end{center}

\paragraph{The auxiliary problem.} Given $h$ and its Fourier transform, determine every real perturbation $k$ for which both $h\pm\varepsilon k$ stay in the cone for some $\varepsilon>0$. A zero on a real frequency line forces divisibility. A nonreal zero of a radial transform instead supplies a polynomial divisor bounded away from zero on real space. We estimate this quotient on shifted contours to control the perturbation at every spatial point. This yields the complete face calculation in Section~\ref{sec:faces}. The cone is taken over real scalars throughout. Section~\ref{sec:xi-face} applies this argument to the classical $\Xi$-function and the de Bruijn--Newman deformation.

Corpus entries 103 and 126 study zeros of entire functions under imaginary translations and backward heat evolution. Entry 027 works with the ball overlaps that appear here as autocorrelations, and entry 029 treats recovery from mixtures and spectral observations \cite{Corpus}. We use these themes to select the auxiliary problem. The estimates below are proved independently. Classical entire-function factorization, sector bounds, Bochner's theorem and the one-dimensional Boas--Kac factorization retain their usual roles.

\section{The extremal-ray questions of Jaming, Matolcsi and R\'ev\'esz}
\label{sec:jmr}
Let $\Omega_d^+$ denote the cone of continuous nonnegative radial positive
definite functions on $\mathbb R^d$. An element generates an extreme ray
if every decomposition in the cone has proportional summands.
Jaming, Matolcsi and R\'ev\'esz~\cite[Questions 1 and 3]{JMR2009}
asked whether products and convolutions preserve extremality. They also
asked specifically whether, for the tent $T(x)=(1-|x|)_+$ and the Gaussian
$\gamma(x)=e^{-\pi x^2}$, $T*\gamma$ generates an extreme ray and, if it
does, whether $\gamma(T*\gamma)$ generates one. Their compact-support
theorem gives extremality when the radial Fourier transform has only real
zeros. We prove that a Gaussian smoothing of every such function is still
extreme, even against unrestricted positive definite, nonnegative
summands. We then exhibit explicit decompositions answering the remaining
questions.

Use $\widehat f(\xi)=\int f(x)e^{-2\pi i x\cdot\xi}\,dx$ and put
$G_\sigma(x)=\sigma^{-d/2}e^{-\pi|x|^2/\sigma}$, $\sigma>0$.
\begin{theorem}[Gaussian smoothing]\label{thm:smoothing}
Let $f$ be a continuous nonzero compactly supported nonnegative positive
definite function on $\mathbb R^d$. For every real unit vector $e$, assume
that the entire function $P_e(z)=\widehat f(ze)$ has only real zeros.
Then $G_\sigma*f$ generates an extreme ray in the cone of all continuous
nonnegative positive definite functions on $\mathbb R^d$. If $f$ is
radial it suffices to test one direction.
\end{theorem}

\begin{proof}
Set $h=G_\sigma*f=g+g_1$ with $g,g_1$ in the cone. Since
$0\le g\le h$ and $h$ has Gaussian decay up to a linear exponential
factor, $g$ is integrable and $\widehat g(ze)$ is entire of order at
most two. By positive definiteness on the two summands,
$0\le\widehat g(\xi)\le\widehat h(\xi)$ on $\mathbb R^d$.
Fix $e$ and put $P=P_e$. Evenness and reality follow from the nonnegative
real-valued positive definite $f$; the same holds for $g$. In particular
$P(0)=\int f>0$. Compact support gives finite exponential type. Pairing
the real zeros in Hadamard's factorization, repeated with their
multiplicities, gives
\[
 P(z)=P(0)\prod_{a>0}(1-z^2/a^2),\qquad \sum_{a>0}a^{-2}<\infty.
\]
The remaining exponential linear factor is constant by evenness.
Since $|\widehat g(xe)|\le e^{-\pi\sigma x^2}P(x)$ for real $x$,
the numerator vanishes to at least the multiplicity of every zero in
this product. Consequently
\[
 L(z)=\frac{\widehat g(ze)e^{\pi\sigma z^2}}{P(z)}
\]
is entire with $|L(x)|\le1$ on the real axis. At $z=x+iy$ spatial
nonnegativity gives
\[
 |\widehat g(ze)|\le\widehat h(iye)=e^{\pi\sigma y^2}P(iy).
\]
When $|x|\le |y|/2$, the ratio of the absolute values of corresponding
factors of $P(iy)$ and $P(z)$ satisfies
\[
 \frac{a^2+y^2}{|a^2-(x+iy)^2|}
 \le \left(1-\frac{x^2}{a^2+y^2}\right)^{-1}
 \le \exp(2x^2/a^2).
\]
Hence $|L(x+iy)|\le e^{D_e x^2}$ on that cone for a finite constant $D_e$.

For completeness, $\log^+|L(z)|=O(1+|z|^2)$ globally.
The corresponding estimate holds for the numerator by spatial decay.
Jensen's formula for $P$ gives the mean of $\log|P|$ on the circle
of radius $R$ at least $\log P(0)$. As $P$ is of exponential type,
the mean of $\log^-|P|$ is $O(1+R)$. It follows that the mean of
$\log^+|L|$ is $O(1+R^2)$. The Poisson bound for the subharmonic
function $\log^+|L|$ on the concentric half-disc gives the assertion.

Fix $\varepsilon>0$. On a ray at angle $\theta<\pi/2$ sufficiently
close to $\pi/2$, we have $\cot\theta\le1/2$ and
$D_e\cos^2\theta\le\varepsilon\sin(2\theta)$. Thus
$e^{i\varepsilon z^2}L(z)$ has modulus at most one both on that ray
and on the positive real axis. The sector Phragm\'en--Lindel\"of
theorem applies because its growth order is at most two and
$2<\pi/\theta$. Let $\varepsilon\downarrow0$ and choose the rays
tending to the positive imaginary axis. The whole first quadrant has
$|L|\le1$. Evenness and conjugation symmetry give the same conclusion
everywhere. Liouville's theorem gives $L(z)=L(0)$. The value at the
origin is $\widehat g(0)/\widehat f(0)$ for every $e$. Uniqueness of
Fourier transforms now gives $g=L(0)h$ throughout $\mathbb R^d$.
\end{proof}

The one-dimensional tent has transform
$\operatorname{sinc}^2(\xi)$, where
$\operatorname{sinc}(u)=\sin(\pi u)/(\pi u)$ with its continuous
extension at zero. Thus $T*\gamma$ is extreme. More generally the
theorem applies to convolutions of tents of any finitely many lengths
and to tensor products of these functions and their invertible linear
images. It also applies directly to the compact radial kernels
$m_\alpha*m_\alpha$ in~\cite[Corollary 3.3]{JMR2009} whenever the
source's hypotheses give only real zeros.

\paragraph{A decomposition after Gaussian tapering.}
For $a,b>0$ let $T_a(x)=(1-|x|/a)_+$ and
$F(x)=e^{-\pi b x^2}T_a(x)$. Its Fourier transform is the convolution
of a positive Gaussian with $a\operatorname{sinc}^2(a\,\cdot)$ and
satisfies
\[
 \widehat F(\xi)\ge\frac{c_{a,b}}{1+\xi^2}.
\]
Indeed, restrict the integration variable $\eta$ to
$[\xi,\xi+1/a]$. The Gaussian factor at $\xi-\eta$ is bounded below;
the denominator in the squared sinc is at most a constant times
$1+\xi^2$, while the integral of $\sin^2(\pi a\eta)$ there equals
$1/(2a)$. If $0<h<a$, $T_h\le C F$ on the real line, since $F$
has a positive minimum on $[-h,h]$, and
$\widehat T_h(\xi)\le C'/(1+\xi^2)$. For sufficiently small
$\delta>0$, both $F-\delta T_h$ and $\delta T_h$ are nonnegative with
nonnegative integrable Fourier transforms. They are nonproportional by
their supports. This proves that every such $F$ is nonextreme.

In particular $T\gamma$ is a product of two extreme generators that is
nonextreme. Fourier transformation gives the same conclusion for the
convolution $\operatorname{sinc}^2*\gamma$. Indeed Fourier
transformation preserves extremality for these integrable cone elements:
each summand is dominated by the given integrable function, and its
Fourier transform is nonnegative and integrable by Bochner's theorem.
Inversion interchanges decompositions. The tent and Gaussian
extremality used here is also established in~\cite{JMR2009}.

Finally, completing the square under the convolution gives
\[
 \gamma(x)(T*\gamma)(x)
 =2\left[e^{-2\pi(\cdot)^2}*
       \left(e^{-2\pi(\cdot)^2}T_{1/2}\right)\right](x).
\]
The inner tapered tent has the displayed decomposition; convolution
with the outer Gaussian preserves positivity and positive definiteness
and is injective on integrable functions, since its Fourier multiplier
is strictly positive. Its two output summands remain nonproportional.
Thus Question~1 has negative answers for both operations, and the two
specified functions in Question~3 give, respectively, an extreme and a
nonextreme ray generator.

\subsection{Counterexamples to both closure assertions in every dimension}
For $a>0$ let $K_a=\mathbf1_{B_a}*\mathbf1_{B_a}$ on $\mathbb R^d$, where the ball is centered at zero. Then $G_1K_a$ and $G_1*\widehat K_a$ each fail to generate an extreme ray of the radial cone. Each input to either operation is extreme. This gives the negative answers to Question~1 separately for every $d\ge1$.

Indeed $K_a$ is extreme by \cite[Cor.~3.3]{JMR2009} and scaling, and duality as above gives the extremality of $\widehat K_a$. The Gaussian is extreme by \cite[Sec.~4.2]{JMR2009}, or by the sector proof above with the constant zero-free profile $P=1$. For $r=|\xi|\to\infty$ the Fourier transform of the ball and the fixed-order Bessel asymptotic give
\[
 \widehat K_a(\xi)=c_{a,d}r^{-(d+1)}
 \left(\cos^2(2\pi a r-(d+1)\pi/4)+O_{a,d}(r^{-1})\right),
 \qquad c_{a,d}>0.
\]
Here one squares a constant multiple of $r^{-d/2}J_{d/2}(2\pi a r)$, using the usual cosine asymptotic \cite[Sec.~10.17]{DLMFBessel}. In particular $\widehat K_b\le C_b(1+|\xi|)^{-(d+1)}$ for every $b>0$. To lower-bound $G_1*\widehat K_a$ at $re$, integrate $\widehat K_a$ on points $(r+v)e+w$, where $0\le v\le 1/a$ and $w\perp e$, $|w|\le1$. Omit the transverse integration in dimension one. The Gaussian is bounded below on this fixed cylinder of displacements. Its frequency radii are $r+v+O(1/r)$ uniformly, and the squared cosine has integral $1/(2a)$ on the longitudinal interval up to an error $O(1/r)$. Thus the convolution is at least $c(1+r)^{-(d+1)}$ for large $r$. It is strictly positive on bounded frequency balls, so the estimate holds everywhere after reducing $c$.

Choose $0<b<a$. The compact support of $K_b$ lies in the open ball where $G_1K_a$ is positive. Hence $K_b\le C G_1K_a$ spatially, and the two spectral estimates give $\widehat K_b\le C'\widehat{G_1K_a}$. Subtracting a sufficiently small positive multiple of $K_b$ gives a nonproportional decomposition in the radial cone. Fourier transformation takes it to a decomposition of $G_1*\widehat K_a$. The same proof works for any positive Gaussian width.

\subsection{Every finite sequence of the two Gaussian operations}
For $s,b>0$ write $H_s f=G_s*f$ and $M_b f(x)=e^{-\pi b x^2}f(x)$ on the line. In any finite sequence of $H_s$ and $M_b$ applied to $T_a$, the output generates an extreme ray exactly when the sequence contains no $M_b$.

Indeed, consecutive heat operators combine by addition of their parameters. A sequence consisting only of these operators is covered by the smoothing theorem and the extremality of the tent. If the first multiplication follows $H_s$, put $v=1+bs$. Completing the square gives
\[
 (M_b H_s f)(x)=\sqrt v\,H_{s/v}\bigl[y\mapsto e^{-\pi b v y^2} f(vy)\bigr](x).
\]
For $f=T_a$ the bracket is a Gaussian times $T_{a/v}$ and has the nonproportional decomposition above. A multiplication applied immediately to the original tent has the same property. All subsequent operations are injective linear maps of the integrable cone into itself, so every such pair of summands stays nonproportional. This proves the classification for arbitrary positive parameters and arbitrary finite lengths.

\subsection{The converse for compact radial cores}
Let $f$ be nonzero, continuous, compactly supported, radial, nonnegative
and positive definite on $\mathbb R^d$. Then $G_\sigma*f$ generates an
extreme ray, for one (equivalently every) $\sigma>0$, if and only if
the even entire function $z\mapsto\widehat f(ze)$ has only real zeros.
The same equivalence holds in the unrestricted cone of continuous
nonnegative positive definite functions. In particular, if $d=1$ the
classification applies to every such compactly supported $f$.

Here is a polynomial-division estimate for Laplace transforms which
also treats some nonradial cores. Suppose $f\ge0$ is integrable with
support in the real ball $B_R$, $F=\widehat f$, and $D$ is a nonconstant
polynomial in $d$ variables such that $U=F/D$ is entire. There is a
constant independent of $z\in\mathbb C^d$ such that
\begin{equation}
 |U(z)|\le C F(i\operatorname{Im}z).
 \label{eq:jmr-polynomial-division}
\end{equation}
To see this, choose a real unit vector $v$ on which the highest homogeneous
part $D_m$ of $D$ is nonzero. For fixed $z$, the polynomial in $w$ given
by $D(z+wv)$ has degree $m$ and leading coefficient $D_m(v)$. Remove
from $[1,2]$ the intervals of radius $1/(4m+4)$ about the absolute values
of its roots. Their combined length is less than one. On the circle
of some remaining radius, the polynomial has modulus at least
$|D_m(v)|(4m+4)^{-m}$. On this same circle,
\[
 |F(z+wv)|\le e^{4\pi R}F(i\operatorname{Im}z)
\]
by nonnegativity and compact support. The maximum modulus principle
on the circle for $w\mapsto U(z+wv)$ proves
\eqref{eq:jmr-polynomial-division}.

Now assume in addition that $f$ is positive definite and that $D$ is even, real and bounded below by a positive constant
on $\mathbb R^d$. With $h=G_\sigma*f$, define $k$ by
$\widehat k(\xi)=e^{-\pi\sigma|\xi|^2}U(\xi)$. Then
$|\widehat k|\le C\widehat h$. Inverting and shifting the contour to
$\mathbb R^d+i x/\sigma$, with the complex Gaussian extended as
$e^{-\pi\sigma z\cdot z}$, gives
\[
 |k(x)|\le C' e^{-\pi|x|^2/\sigma} F(ix/\sigma).
\]
Indeed, completing the square leaves the real integrable Gaussian
$e^{-\pi\sigma|\xi|^2}$ times $U(\xi+ix/\sigma)$, bounded by
\eqref{eq:jmr-polynomial-division}. The contour shift follows
coordinate by coordinate: on each fixed strip \eqref{eq:jmr-polynomial-division}
bounds $U$ uniformly in all real parts, and the Gaussian sends the
vertical boundary contributions to zero. Directly,
\[
 h(x)=\sigma^{-d/2}e^{-\pi|x|^2/\sigma}
   \int f(y)e^{2\pi x\cdot y/\sigma}e^{-\pi|y|^2/\sigma}\,dy
 \ge \sigma^{-d/2} e^{-\pi R^2/\sigma}
       e^{-\pi|x|^2/\sigma}F(ix/\sigma).
\]
Consequently $|k|\le C''h$. Evenness and reality make $k$ real.
For small $\varepsilon>0$, $h+\varepsilon k$ and $h-\varepsilon k$
are nonnegative with nonnegative integrable Fourier transforms. They
are nonproportional, since $D$ is nonconstant and $F$ is nonzero
in a neighborhood of zero.

It remains to obtain this divisor in the radial case. Write
$P(s)=F(se)$ and $P(s)=A(s^2)$ with $A$ entire.
The equality $F(z)=A(z\cdot z)$ holds on $\mathbb R^d$ by radiality
and hence on $\mathbb C^d$ by analytic continuation. If $\zeta$
is a nonreal zero of $P$, then $w=\zeta^2$ is either negative real
or nonreal. In the former case take $d_0(v)=v-w$, and in the latter
take $d_0(v)=(v-w)(v-\overline w)$. In both cases $A/d_0$ is entire,
and $D(z)=d_0(z\cdot z)$ is a nonconstant real even polynomial with
a positive lower bound on $\mathbb R^d$. The two perturbations just
constructed are radial, so also refute extremality in the smaller
radial cone. The converse implication is the Gaussian smoothing theorem.

A compact nonzero continuous $f$ as above has a Fourier transform
along each ray tending to zero at real infinity. If $P$ has no real
zeros but were zero-free, Hadamard factorization and evenness would make
it constant. Consequently, if a Gaussian smoothing of a nonzero radial compact core and its Fourier transform are both strictly positive everywhere, the smoothing cannot generate an extreme ray.

\section{The dimension of the face after heat flow}
\label{sec:faces}
The proof of Theorem~3.2 in \cite{JMR2009} bounds the degree of a polynomial quotient for unsmoothed compact radial functions. Here the Gaussian contour estimate realizes every even polynomial up to that degree. For an integrable member $h$ of the cone of all continuous nonnegative positive definite functions, let $\mathcal V(h)$ be the real vector space of real even integrable functions $k$ with continuous representatives satisfying
\[
 |k|\le C h,\qquad |\widehat k|\le C\widehat h
\]
for some finite $C$. For every such $k$, the two functions $h\pm\varepsilon k$ belong to the cone for all sufficiently small $\varepsilon>0$. Conversely, two-sided perturbations in the cone satisfy these inequalities by integrability and Bochner's theorem. Thus $\mathcal V(h)$ is the linear span of the smallest face containing $h$. Indeed that face consists of the cone elements $g$ for which $ch-g$ is in the cone for some $c>0$: this set is a face by addition and the cone order, and each such $g$ lies in every face containing $h$ by the decomposition $ch=g+(ch-g)$. In the radial cone take the radial subspace of $\mathcal V(h)$.

\begin{theorem}\label{thm:faces}
Let $h=G_\sigma*f$ with $f$ as in the radial converse above. Count the nonreal zeros of $P(t)=\widehat f(te)$ with multiplicities. If there are infinitely many, the radial face span is infinite dimensional. If their number is $m<\infty$, write $D(t)=d_0(t^2)$ for the monic polynomial collecting them, with $D=1$ when $m=0$, and $P_0=P/D$. In the full cone the perturbations have exactly the Fourier transforms
\[
 \widehat k(\xi)=e^{-\pi\sigma|\xi|^2} P_0(|\xi|) Q(\xi),
 \qquad Q\in\mathbb R[\xi_1,\ldots,\xi_d],\quad
 Q(-\xi)=Q(\xi),\quad \deg Q\le m.
\]
The full and radial face dimensions are respectively
\[
 \sum_{j=0}^{m/2}\binom{d+2j-1}{2j}
 \quad\text{and}\quad 1+m/2.
\]
Each finite radial dimension is odd, and every positive odd integer occurs in each dimension.
\end{theorem}

\begin{proof}
$P(iy)=\int f(x)e^{2\pi y x\cdot e}\,dx>0$ for every $y$. Thus every nonreal zero is in a quartet under negation and conjugation, and $m$ is a multiple of four. In particular $D$ and $d_0$ are positive on their respective real and nonnegative real arguments. The even entire $P_0$ has exponential type by polynomial division on the line and has only real zeros.

Fix $k\in\mathcal V(h)$ and a real unit vector $e$. Its Gaussian spatial bound makes its line transform entire. Set $W(t)=\widehat k(te)e^{\pi\sigma t^2}/P_0(t)$. The real zeros divide the numerator with multiplicity, by the spectral bound. On the real line $|W(t)|\le C(1+|t|)^m$. The Laplace estimate in Section~\ref{sec:jmr} applies to $|k|\le Ch$. Its real-zero product argument applied to $P_0$ gives
\[
 |W(x+iy)|\le C(1+|x+iy|)^m e^{A_e x^2}
 \qquad(|x|\le |y|/2).
\]
Indeed the extra factor is $|D(iy)|$. Jensen's estimate again gives $\log^+|W(t)|=O(1+|t|^2)$. In the first quadrant, $|t+i|\ge (1+|t|)/\sqrt2$. Apply the same subcritical sector argument to $e^{i\varepsilon t^2}W(t)/(t+i)^m$. The boundary constant is uniform as the second ray tends to the imaginary axis. Letting $\varepsilon\downarrow0$ gives a polynomial bound of degree $m$ throughout the quadrant. Reality and evenness extend it to the plane. Cauchy's estimates now give a polynomial $W$ of degree at most $m$.

Near $z=0$ the entire numerator in $d$ variables divided by the lifted $P_0(\sqrt{z\cdot z})$ is holomorphic since $P_0(0)>0$. Its Taylor polynomial of degree $m$ agrees with $W$ identically on each real line through zero, because it agrees there through degree $m$. This supplies one real even multivariate polynomial $Q$ as asserted.

Conversely let any such $Q$ be given. Its spectral quotient $Q(\xi)/D(|\xi|)$ is bounded. For the spatial bound use $U(z)=F(z)/d_0(z\cdot z)$ in \eqref{eq:jmr-polynomial-division}, or $U=F$ if $m=0$, where $F=\widehat f$. On the shifted contour $z=t+iy$, where $y=x/\sigma$, the absolute inverse transform of the proposed spectrum is bounded by
\[
 e^{-\pi|x|^2/\sigma}\int_{\mathbb R^d}e^{-\pi\sigma |t|^2}|U(t+iy)Q(t+iy)|\,dt.
\]
For large $|y|$ and $|t|\le |y|/2$, the real part of $(t+iy)\cdot(t+iy)$ has absolute value at least $3|y|^2/4$. Hence $|d_0((t+iy)\cdot(t+iy))|\ge c |y|^m$, while $|Q(t+iy)|\le C(1+|y|)^m$. This bounds $|UQ|$ by $C F(iy)$ there. On $|t|>|y|/2$, the division estimate bounds $|UQ|$ by $C F(iy)(1+3|t|)^m$, integrable against the Gaussian with uniform constant. Bounded $|y|$ is covered directly by the division estimate. As in the preceding section, the shifted bound gives $|k|\le C'h$ and justifies the contour move. This proves the characterization. Counting even monomials gives the stated dimension. In the radial subspace $Q$ is a polynomial in $|\xi|^2$, as is seen by restricting to the first axis and using rotation invariance.

If there are infinitely many nonreal zeros there are infinitely many distinct quartets, since each zero has finite multiplicity. Divide $F$ by one quartet polynomial at a time as in \eqref{eq:jmr-polynomial-division}. The resulting radial perturbations are linearly independent by their distinct poles after division of their transforms by $\widehat h$.

To realize dimension $1+2k$ for $k\ge1$, fix an integer $\alpha>4k$ and put $u(x)=(1-|x|^2)_+^\alpha$, $v=u+\delta\Delta^{2k}u$. The cutoff $u$ is $C^{4k}$, so the derivatives are ordinary continuous functions. With $r=|x|$ approaching $1$ from below, the leading term of $\Delta^{2k}u$ is
\[
 (2r)^{4k}\alpha(\alpha-1)\cdots(\alpha-4k+1)(1-r^2)^{\alpha-4k}.
\]
All remaining terms are $O((1-r^2)^{\alpha-4k+1})$ by the chain rule. Hence $v\ge0$ near the boundary, and a sufficiently small $\delta>0$ gives $v\ge0$ on the remaining compact ball since $u>0$ there. Set $f=u*v$. The line transform of $u$ is a normalized Bessel function times a positive constant and has only real zeros \cite[Cor.~3.3]{JMR2009}. The spectral density of $f$ is $|\widehat u(\xi)|^2(1+\delta(2\pi |\xi|)^{4k})\ge0$. The polynomial adds precisely $4k$ nonreal zeros on the line and the physical convolution is nonnegative. Taking $f=u*u$ realizes dimension one.
\end{proof}

\section{The minimal face of the square of the Riemann Xi function}
\label{sec:xi-face}
Write $\xi(s)=s(s-1)\pi^{-s/2}\Gamma(s/2)\zeta(s)/2$ and
$\Xi(z)=\xi(1/2+iz)$. The function $\Xi^2$ itself belongs to
$\Omega_1^+$. Its minimal face has real dimension $1+N$, where $N$
is the number of nontrivial zeta zeros off the critical line, counted
with multiplicity, with the convention that infinitely many such zeros
give infinite dimension. In particular, the Riemann hypothesis is
equivalent to the assertion that $\Xi^2$ generates an extreme ray.

We prove the statement throughout the de Bruijn--Newman family. Use
the normalization of Rodgers and Tao~\cite{RTNewman}:
\[
 \Phi(v)=\sum_{n\ge1}
  (2\pi^2n^4e^{9v}-3\pi n^2e^{5v})e^{-\pi n^2e^{4v}},\qquad
 H_t(z)=\int_0^\infty e^{tv^2}\Phi(v)\cos(zv)\,dv.
\]
The theta transformation makes $\Phi$ even, and every summand is positive
for $v\ge0$. In particular
\[
 H_0(z)=\tfrac18\Xi(z/2),\qquad
 u_t(x)=\pi e^{4\pi^2tx^2}\Phi(2\pi x)>0,\qquad
 \widehat u_t(z)=H_t(z).
\]
The transform identity uses the Fourier convention with exponent
$-2\pi i xz$. If $u$ is the inverse Fourier transform of $\Xi$ itself,
then $u(x)=4u_0(x/2)$. For fixed $t\in\mathbb R$ and $\sigma\ge0$, set
\[
 h_{t,\sigma}=G_\sigma*(u_t*u_t),\qquad
 \widehat h_{t,\sigma}(x)=e^{-\pi\sigma x^2}H_t(x)^2,
\]
with $G_0$ denoting the unit mass at zero. All these functions and their Fourier transforms are integrable, nonnegative and positive definite.
Indeed, the spatial functions are nonnegative by convolution, their
transforms are nonnegative, and $H_t^2$ is integrable by Plancherel
since $u_t\in L^2$. Write $N_t$ for the number of nonreal zeros of
$H_t$, counted in the whole complex plane with multiplicity.

\paragraph{Face-dimension formula.}
The dimensions of the minimal faces of $h_{t,\sigma}$ and of
$e^{-\pi\sigma x^2}H_t(x)^2$ are both $1+N_t$ when $N_t<\infty$,
and both are infinite otherwise. In particular $H_t^2$ generates an
extreme ray precisely for $t\ge\Lambda$. For every $t<0$ the face is
infinite-dimensional by Dobner's theorem~\cite{DobnerNewman}.
For every $t>0$ its dimension is finite by Ki, Kim and
Lee~\cite{KKLNewman}. At $t=0$ the formula for $\Xi^2$ follows by
scaling. The same criterion holds for $G_\sigma*(u*u)$ at every
$\sigma\ge0$, by the transform identity and Gaussian scaling.
These equivalences leave the zero count $N_0$ undetermined.

To prove the formula, fix $t,\sigma$ and abbreviate $P=H_t$, $u=u_t$,
$h=h_{t,\sigma}$. Two estimates on the theta kernel will be used.
For some $b,c,C>0$, $u(x)\le C e^{-c e^{b|x|}}$. Moreover for
every $A>0$ there is $X_A$ such that
\begin{equation}
 u(x+s)\le e^{-As}u(x)\qquad (x\ge X_A,\ s\ge0).
 \label{eq:xi-tail}
\end{equation}
For the first estimate separate the $n=1$ term and sum the exponentially
decreasing tail for $v\ge0$. For the second write
$w=\pi n^2 e^{4v}$. The logarithmic derivative of each positive
summand in $\Phi(v)$ is $5+8w/(2w-3)-4w$. Uniformly over $n\ge1$
this tends to $-\infty$ at least as fast as $-c'e^{4v}$ as
$v\to\infty$. The additional derivative $2tv$ from $e^{tv^2}$
does not change this conclusion. Integrating the common upper bound
and summing proves \eqref{eq:xi-tail}; evenness gives the left counterpart.
The double-exponential bound yields
$\log^+\max_{|z|\le R}|P(z)|=O_t((R+1)\log(R+2))$: substitution
$v=c e^{b x}$ in the absolute Fourier integral bounds it by a constant
times a rescaled gamma integral. Thus $P$ has order at most one and
$P(iy)>0$ for all $y\in\mathbb R$.

The linear span of the minimal face of $h$ consists of the real even
continuous functions $k$ satisfying
\[
 |k|\le C h,\qquad |\widehat k|\le C\widehat h
\]
for some $C$. Indeed the minimal face is the order ideal
$\{g\in\Omega_1^+: C h-g\in\Omega_1^+\text{ for some }C>0\}$:
it is a face and every face containing $h$ contains it. The two
inequalities are equivalent to $h\pm\varepsilon k\in\Omega_1^+$
for some $\varepsilon>0$ by integrability and Fourier inversion.
Fourier transformation identifies the two face spans, so it is enough
to consider this one.

First suppose $N=N_t<\infty$, and let $D$ be the monic polynomial
whose zeros are the nonreal zeros of $P$ with their multiplicities.
It is real and even, positive on the real axis, with degree $N$; use
$D=1$ when $N=0$. The function $V=P/D$ is even of order at most one.
As in the Gaussian smoothing theorem its Hadamard product has the form
$V(z)=V(0)\prod_{a>0}(1-z^2/a^2)$, with only real zeros and
$\sum a^{-2}<\infty$. For each $k$ in the face span consider
\[
 L(z)=e^{\pi\sigma z^2}\widehat k(z)/V(z)^2.
\]
The bound by $Ch$ makes $\widehat k$ entire, and the real-axis spectral
bound cancels every real zero of $V^2$. Thus $L$ is entire and
$|L(x)|\le C D(x)^2$ for real $x$. The circle-mean argument of the
Gaussian smoothing theorem gives
$\log^+|L(z)|=O(1+|z|^2)$, since
$|\widehat k(x+iy)|\le C e^{\pi\sigma y^2}P(iy)^2$.
In the cone $|x|\le |y|/2$ the real-zero product gives, as before,
\[
 |L(x+iy)|\le C(1+|z|)^{2N}e^{B x^2}
\]
for some $B$. Divide by $(z+i)^{2N}$ in the first quadrant. This
analytic quotient is bounded on the positive real axis, and bounded
by $C'e^{B x^2}$ in the indicated part of the first quadrant.
Applying the same sector argument to the quotient times
$e^{i\varepsilon z^2}$ and then letting $\varepsilon\downarrow0$
gives $|L(z)|\le C''(1+|z|)^{2N}$ in the quadrant. Evenness and
conjugation extend this bound everywhere. Cauchy's inequalities show
that $L$ is an even real polynomial of degree at most $2N$. The face
span therefore has dimension at most $N+1$. This bound includes
$\sigma=0$.

For the reverse inequality, $D$ may now be any finite product of
complete nonreal zero orbits of $P$, including repeated factors up
to their available multiplicities, and write $n=\deg D$.
For $0\le j\le n$ let $v_j$ have Fourier transform
$(i x)^jP(x)/D(x)$. We show $|v_j|\le C_j u$. By partial fractions
it suffices to invert $P(x)/(x-\alpha)^\ell$, where $P$ vanishes at
$\alpha$ to order at least $\ell$. The constant polynomial term
contributes a multiple of $u$. For the proper fractions use
\[
 \frac{(2\pi i)^\ell}{(\ell-1)!}
 \int_{-\infty}^x (x-y)^{\ell-1}
             e^{2\pi i\alpha(x-y)}u(y)\,dy .
\]
The vanishing of $P^{(j)}(\alpha)$ for $j<\ell$ shows that the same
expression equals the negative of the integral over $[x,\infty)$
with the displayed integrand and constant. The right integral for
large positive $x$ and the left one for large negative $x$ have
modulus at most $C u(x)$ by \eqref{eq:xi-tail}, taking
$A>2\pi|\operatorname{Im}\alpha|+1$. On a compact interval, positivity
and continuity of $u$ give the same comparison. Applying
$((2\pi i)^{-1}d/dx-\alpha)^\ell$ distributionally returns $u$,
which verifies the Fourier transform.

Now set $k_j=G_\sigma*(v_j*\widetilde v_j)$ with
$\widetilde v_j(x)=\overline{v_j(-x)}$. The spatial bound just proved
gives $|k_j|\le C_j^2 h$. Its Fourier transform is
\[
 \widehat k_j(x)=e^{-\pi\sigma x^2} x^{2j}P(x)^2/D(x)^2.
\]
Because $x^{2j}/D(x)^2$ is bounded on $\mathbb R$ for $j\le n$,
all $k_j$ lie in the face span. The distinct even monomials give
linear independence on an interval about zero. Taking the full $D$
in the finite case proves equality. In the infinite case take finite
orbit products with arbitrarily large degree.

The classical positivity criterion of Weil and Bombieri's result on the
negative eigenvalues of sufficiently large finite truncations of Weil's
quadratic functional~\cite{BombieriWeil} also relate positivity to
off-line zeros. Csordas~\cite{CsordasXi} treats the positive definiteness
of an associated sequence of kernels via generalized Laguerre inequalities.
Here the invariant is the dimension of the minimal face containing
the single function $H_t^2$ in the positive positive definite cone.

\section{A tail criterion for the Gaussian uniqueness question}
\label{sec:tails}
Borisov asked about other extreme rays among his good functions; Jaming, Matolcsi and R\'ev\'esz formulated the one-dimensional strict-positivity conjecture in the cone $\Omega_1^+$ \cite[Sec.~5]{Borisov}\cite[Sec.~6]{JMR2009}. Here strict positivity means $f(x)>0$ and $\widehat f(\xi)>0$ on their entire real axes. Theorem~\ref{thm:tail} establishes the conjectured conclusion on two classes specified by the physical tail. Fourier symmetry gives the same conclusion when the hypotheses hold on the frequency side. We use the $2\pi$ convention of Section~\ref{sec:jmr}.

\begin{theorem}\label{thm:tail}
Suppose $f\in\Omega_1^+$ is integrable, $f,\widehat f$ are strictly positive, and $V(x)=-\log f(x)$ is convex for sufficiently large positive $x$. Assume either
\begin{enumerate}
\item $V'_+(x)$ has a finite limit as $x\to\infty$, or
\item $V(x)/(x\log x)\to\infty$.
\end{enumerate}
Then $f$ generates an extreme ray exactly when it is a positive multiple of a centered Gaussian.
\end{theorem}

\begin{proof}
A strictly positive continuous integrable function with the prescribed convex tail tends to zero there, and convexity then makes $V$ eventually increasing. In the first case the tail has Lipschitz constant $L<\infty$. Evenness and bounded oscillation of $V$ on the complementary compact interval give
\[
 f(x-y)\le C e^{L|y|} f(x)\qquad(x,y\in\mathbb R).
\]
For example, extend the two tails across the central interval by a constant and absorb the bounded difference from $V$ into $C$. Choose $a>L$ and $k=f*e^{-a|\cdot|}$. Spatially $0<k\le C' f$. Its Fourier transform is $\widehat f(\xi)\,2a/(a^2+4\pi^2\xi^2)$, a bounded nonconstant multiple of the strictly positive density. Thus $f+\varepsilon k$ and $f-\varepsilon k$ belong to the cone for sufficiently small $\varepsilon>0$ and are nonproportional.

In the second case, $F(z)=\widehat f(z)$ is entire. For each $B>0$ there is a constant $C_B$ with $f(x)\le C_B\exp[-B|x|\log(2+|x|)]$. For $|z|\le r$ we may take $e^{-|x|}$ out of the integral and bound the remaining exponential by its supremum. Using $\sup_{t>0}(ut-Bt\log t)=B\exp(u/B-1)$ with $u=2\pi r+1$ gives
\[
 \log^+\max_{|z|\le r}|F(z)|\le C'_B+B\exp((2\pi r+1)/B).
\]
Taking a further positive logarithm, dividing by $r$, and letting $B$ increase yields the growth bound $\log^+\log^+\max_{|z|\le r}|F(z)|=o(r)$. If $F$ has no zeros, $F/F(0)$ is the entire characteristic function of the probability density $f/\int f$. Ostrovskii's theorem, in the formulation of Eremenko and Fryntov \cite{EF}, makes that density Gaussian. It is centered by evenness. Its extremality follows from \cite[Sec.~4.2]{JMR2009}.

Suppose, therefore, that $F$ has a zero $z_0$. All zeros lie off the real axis. Let $D$ be the monic polynomial with the distinct roots in the orbit of $z_0$ under conjugation and negation. Since $F$ is real entire and even, $D$ divides $F$; $D$ is even, real, and strictly positive on the real line. Define $k$ by $\widehat k=F/D$. We verify the bound $|k|\le C f$, including in both tails.

Write $\mathcal L=(2\pi i)^{-1}\,d/dx$. For a root $z_j$ of $D$, divide a rapidly decreasing function $g$ by the factor $\xi-z_j$ on the Fourier side by taking
\[
 u(x)=-2\pi i\int_x^\infty e^{2\pi i z_j(x-t)}g(t)\,dt.
\]
Then $(\mathcal L-z_j)u=g$. If $\widehat g(z_j)=0$, the same expression is $2\pi i\int_{-\infty}^x e^{2\pi i z_j(x-t)}g(t)\,dt$. Start with $g=f$ and iterate over the distinct roots. Convexity and $V(x)/(x\log x)\to\infty$ give $V'_+(x)\to\infty$. Hence for every $A>0$, eventually $f(x+s)\le f(x)e^{-As}$ for $s\ge0$, and evenness gives the corresponding left-tail bound. If $|g|\le C f$, either integral at its respective end is at most $C' f(x)$ by taking $A>2\pi|\operatorname{Im}z_j|+1$. Continuity and strict positivity of $f$ absorb the remaining compact interval. This proves $|u|\le C'' f$ at each step. In particular all transforms remain entire, and the differential equation identifies $\widehat u=\widehat g/(\xi-z_j)$ by integration by parts and continuation. Thus the next root is still a zero as required. After the last step the function is the real even $k$ by its real even Fourier transform. Since $D$ is bounded below on the real line, $|\widehat k|\le C\widehat f$. As in the first case the two perturbations $f\pm\varepsilon k$ disprove extremality. They cannot be proportional because $1/D$ is nonconstant and $\widehat f$ is strictly positive.
\end{proof}

The second condition includes tails with $V(x)\ge c x^{1+\delta}$ for some $c,\delta>0$. The theorem leaves open, among other cases, tails at scale $x\log x$ and positive functions with uncontrolled oscillation in both domains.

\section{Tests for mixtures of ball autocorrelations}
\label{sec:ballmix}
For $d\ge1$, $\alpha>-1/2$, and $s>0$ put
\[
 m_\alpha(x)=(1-|x|^2)^\alpha\,\mathbf1_{\{|x|<1\}},\quad
 K_\alpha=m_\alpha*m_\alpha,\quad
 C_\alpha=K_\alpha(0)
 =\frac{\pi^{d/2}\Gamma(2\alpha+1)}{\Gamma(d/2+2\alpha+1)}.
\]
At the boundary define $m_\alpha=0$. Question 2(a) in
\cite{JMR2009} asks for a test of membership in scale mixtures of these
kernels, with shape fixed or variable. We first use the bounded positive
mixing measures of their Lemma 2.1, retaining the raw kernels
$K_\alpha(x/s)$. For a variable shape we also give the corresponding test
for normalized kernels $K_\alpha/C_\alpha$.

A continuous radial candidate $f$ must be positive definite. This can
itself be tested by the usual finite positive-semidefinite matrices.
Use angular frequencies in this section, so its Bochner measure $\sigma_f$
satisfies $f(x)=\int e^{i x\cdot\xi}\,d\sigma_f(\xi)$.
For strictly positive finite scales the tests below first require $\sigma_f(\{0\})=0$. If the multiplier formulation $K_\alpha(rx)$ admits $r=0$, subtract $c=\sigma_f(\{0\})$ and apply the test to $f-c$. The atom is uniquely determined by the Gaussian mean below. Its constant contribution is supplied by the zero multiplier with mixing mass $c/C_\alpha$ at any one allowed shape; the minimum-mass statement below concerns the positive-scale residual.
Define $\rho_f$ on $\mathbb R$ by pushing this measure forward by
$\xi\mapsto\log|\xi|$. Isotropy and uniqueness in Bochner's theorem show
that $\rho_f$ completely determines $f$. A zero candidate is represented
by the zero mixing measure; otherwise $f(0)>0$.

\paragraph{The spectral multiplier.}
Write $D=d/2$. Define
\begin{equation}\label{eq:ball-Mellin}
 M_\alpha(z)=
 \frac{\Gamma(D+z/2)}{\Gamma(D)}
 \frac{\Gamma(\alpha+1/2-z/2)}{\Gamma(\alpha+1/2)}
 \frac{\Gamma(\alpha+1)}{\Gamma(\alpha+1-z/2)}
 \frac{\Gamma(D+2\alpha+1)}{\Gamma(D+2\alpha+1-z/2)} .
\end{equation}
If $Y_\alpha$ is the log spectral radius for $K_\alpha/C_\alpha$, then
$\mathbb E e^{zY_\alpha}=M_\alpha(z)$ for $-d<\Re z<2\alpha+1$.
In particular, $M_\alpha(it)$ has no real-frequency zeros.
Indeed, the Fourier--Bessel formula for $m_\alpha$
\cite[Corollary 3.3 and its proof]{JMR2009} gives the ordinary
spectral-radius density proportional to
$r^{-2\alpha-1}J_{D+\alpha}(r)^2$. Insert
$\lambda=2\alpha+1-z$, $\nu=D+\alpha$ into the equal-order
Weber--Schafheitlin integral
\[
 \int_0^\infty J_\nu(r)^2 r^{-\lambda}\,dr
 =\frac{\Gamma(\lambda)\Gamma(\nu+(1-\lambda)/2)}
 {2^\lambda\Gamma((\lambda+1)/2)^2\Gamma(\nu+(1+\lambda)/2)}
 \quad (0<\Re\lambda<2\nu+1).
\]
This is \href{https://dlmf.nist.gov/10.22.E57}{DLMF (10.22.57)}.
Divide by the value at $z=0$ and apply the duplication formula.

\paragraph{Fixed shape.}
For any prescribed $\alpha>-1/2$, the necessary and sufficient
remaining condition is that the following function be positive definite
on the additive group $\mathbb R$:
\begin{equation}\label{eq:ball-divide}
 t\longmapsto\frac{\int e^{ity}\,d\rho_f(y)}{f(0)M_\alpha(it)}.
\end{equation}
Its value at zero is one and it is continuous. In fact, the probability
law with this characteristic function is exactly that of $-\log s$ under
the normalized scale measure. Scaling $x\mapsto x/s$ translates the
log spectral law by $-\log s$. This proves necessity. Conversely,
Bochner's theorem on $\mathbb R$ supplies the translation law;
convolving it with the distribution of $Y_\alpha$ gives
$\rho_f/f(0)$ by Fourier uniqueness. Isotropy recovers $f$,
and its raw mixing mass is $f(0)/C_\alpha$. This proves sufficiency and
uniqueness for arbitrary finite positive scale measures, including
atomic and heavy-tailed ones.

\paragraph{Variable shape, bounded raw mixing measure.}
For a real $h\in C_0(\mathbb R)$ put
\[
 V_h(\alpha,u)=\frac{C_\alpha}{1+C_\alpha}
                   \mathbb E h(Y_\alpha+u),\qquad
 L(h)=\max\{0,\sup_{\alpha>-1/2,\ u\in\mathbb R}V_h(\alpha,u)\}.
\]
Then $f(x)=\int K_\alpha(x/s)\,d\mu(\alpha,s)$ for a finite
positive Borel measure $\mu$ if and only if there is a finite $H\ge0$
such that
\begin{equation}\label{eq:ball-all}
             \int h\,d\rho_f\ \le\ H L(h)
                   \qquad\text{for every real }h\in C_0(\mathbb R).
\end{equation}
The same condition results from restricting $H$ to nonnegative
integers, and $h$ to finite rational linear combinations of the functions
$h_q(y)=e^{-(y-q)^2}$, $q\in\mathbb Q$. For each such test the
two-parameter supremum uses only the explicit integrals
\begin{equation}\label{eq:ball-Gauss-test}
 \mathbb E h_q(Y_\alpha+u)
  =\frac1{2\sqrt\pi}
    \int_{\mathbb R} e^{-t^2/4}
         \operatorname{Re}\{e^{it(u-q)}M_\alpha(it)\}\,dt .
\end{equation}
Both quantities on the input side can be evaluated directly from $f$.
Extend $a_q(\xi)=e^{-(\log|\xi|-q)^2}$ by zero at the origin.
This is a Schwartz function, since every derivative is bounded at zero
and infinity by an arbitrarily high power of the respective decay
parameter. With angular-frequency transform
$\mathcal Fa_q(x)=\int e^{-ix\cdot\xi}a_q(\xi)\,d\xi$,
Fourier inversion and Fubini give
\[
 \int h_q\,d\rho_f=(2\pi)^{-d}\int_{\mathbb R^d}
             f(x)\mathcal Fa_q(x)\,dx,\qquad
 \sigma_f(\{0\})=\lim_{R\to\infty}
 \frac{\int f(x)e^{-|x|^2/(2R^2)}\,dx}{(2\pi R^2)^{d/2}}.
\]
The second identity follows by dominated convergence on the spectral
side. Allowing real $H$, the smallest feasible $H$ equals
$f(0)+\min\mu(({-}1/2,\infty)\times(0,\infty))$, the minimum over
representations of $f$.

\begin{proof}
Let $P=(-1/2,\infty)\times\mathbb R$ with coordinates $(\alpha,u)$,
and set $\tau_{\alpha,u}=(C_\alpha/(1+C_\alpha))
\operatorname{Law}(Y_\alpha+u)$. Formula \eqref{eq:ball-Mellin} and
the characteristic-function continuity theorem show weak continuity
for parameters in $P$. We claim the map extends continuously to the
one-point compactification $P^*$ by assigning the zero measure at infinity,
in the weak-star topology of $C_0(\mathbb R)^*$. If $\alpha\to\infty$,
$C_\alpha\to0$ by the gamma-ratio asymptotic. As
$\alpha\downarrow-1/2$, \eqref{eq:ball-Mellin} gives
$M_\alpha(it)\to0$ for each $t\ne0$. Any bounded interval indicator is
bounded by a constant times a translate of $e^{-y^2}$. Gaussian Fourier
inversion bounds its mass under \emph{every} translation of $Y_\alpha$
by a fixed multiple of
$\int e^{-t^2/4}|M_\alpha(it)|\,dt$, which tends to zero by dominated
convergence. Finally, when $\alpha$ ranges over a compact subinterval,
the weakly continuous family of laws of $Y_\alpha$ is uniformly tight,
so translations with $|u|\to\infty$ have vanishing mass on every compact.
These three cases prove the claim.

The barycenters of the family $\tau_p$, $p\in P^*$, with probability
laws on $P^*$ form a compact convex subset of $C_0(\mathbb R)^*$ in
that topology. Its support function is $L(h)$. Hahn--Banach separation
therefore gives, for $H>0$, membership of $\rho_f/H$ precisely under
\eqref{eq:ball-all}. Dropping the atom at the added point gives a positive
measure $\eta$ on $P$ of mass at most $H$ representing $\rho_f$ by
the $\tau_{\alpha,u}$. Taking $s=e^{-u}$ and
$d\mu=d\eta/(1+C_\alpha)$ produces the original raw representation.
Conversely any such $\mu$ gives a representing
$d\eta=(1+C_\alpha)d\mu$ of mass $\mu(P)+f(0)$ after this change of
coordinates. Taking a decreasing sequence of feasible real $H$ in the
inequalities shows their infimum is feasible. The preceding formula
then proves the minimum assertion (the case $H=0$ gives $f=0$).
The real span of the rational translates of $e^{-y^2}$ is dense in
$C_0(\mathbb R)$: a finite signed measure annihilating that span
has continuous Gaussian convolution zero on $\mathbb Q$, hence zero
everywhere, and the Gaussian's nonvanishing Fourier transform implies
the measure is zero. Rational coefficients suffice by approximation.
\end{proof}

\paragraph{A spatial test that also certifies positive definiteness.}
For a bounded, continuous, radial, real candidate $f$, no preliminary
Bochner test is needed if one uses $h\in C_c^\infty(\mathbb R)$.
Require the displayed Gaussian spatial average to tend to zero and
replace the left side of \eqref{eq:ball-all} by
\[
 T_f(h)=(2\pi)^{-d}\int f(x)\,
    \mathcal F\big[h(\log|\cdot|)\big](x)\,dx.
\]
The radial multiplier is smooth and supported in a compact annulus.
Necessity follows by pairing with a spectral measure.
Conversely the inequalities for $h$ and $-h$ give
$|T_f(h)|\le H\|h\|_\infty$. This functional extends to a real
bounded functional on $C_0(\mathbb R)$ and still satisfies
\eqref{eq:ball-all}, since $|L(h)-L(g)|\le\|h-g\|_\infty$.
The separation proof gives a raw mixture $g$ with the same functional.
The Fourier distribution of $f-g$ pairs to zero against radial
test functions supported away from the origin. Averaging every test
function over rotations gives the same pairing by radiality, so its
support is contained in $\{0\}$. A distribution supported at one point
is a finite sum of derivatives of its Dirac mass; hence $f-g$ is a
polynomial. Boundedness makes it constant and the imposed Gaussian
average makes the constant zero. Thus the inequalities alone certify
both the representation and positive definiteness, with boundedness
and the zero-mean limit as necessary preliminary checks.

\paragraph{Variable shape after peak normalization.}
If the finite measure is instead placed on the normalized profiles
$k_\alpha(r)=K_\alpha(r e_1)/C_\alpha$, an exact test on spatial
values is as follows. Normalize $f(0)=1$ and write $F(r)=f(r e_1)$.
Let $E_m=[-1/2+1/m,m]\times[1/m,m]$ for $m\ge 2$.
There must exist nondecreasing integers $m_j\ge2$ such that every finite
choice of rational $r_i\ge0$ and $b_i\in\mathbb Q$ satisfies
\[
 \sum_i b_i F(r_i)\ \le \sum_{j\ge1}2^{-j}
 \max_{(\alpha,s)\in E_{m_j}}\sum_i b_i k_\alpha(r_i/s).
\]
This condition is sufficient as well. On compact subsets of the shape
range $m_\alpha$ depends continuously on $\alpha$ in $L^2$, by dominated
convergence. Thus each maximum exists. The tail beyond $J$ is bounded
in absolute value by $2^{-J}\sum_i|b_i|$. To prove necessity, split a
representing probability iteratively into pieces of masses $2^{-j}$,
each supported on some $E_{m_j}$: the remaining mass before step $j$ is
$2^{1-j}$ and some compact in the exhaustion contains more than half of it;
take a proportional restriction. For sufficiency, take the product of
the compact spaces of probabilities on these $E_{m_j}$. The coordinate
map taking their $2^{-j}$-weighted integrals at rational radii is
continuous, with compact convex image in the countable coordinate
product. Its support functions are exactly the displayed right sides.
Finite-coordinate separation, rational approximation of coefficients,
and continuity in the radius give the result. The exhaustion matters:
$k_\alpha(r/\sqrt\alpha)\to e^{-r^2/2}$, since
$m_\alpha(\cdot/\sqrt\alpha)\to e^{-|\cdot|^2}$ in $L^2$.
This Gaussian is extreme by \cite[Sec.~4.2]{JMR2009}.
An exact normalized barycenter of an extreme point is concentrated on
that point: splitting at level sets of any rational evaluation proves
this simultaneously at all rational radii. Each $k_\alpha$ has finite
support, so the Gaussian is excluded.

\paragraph{Boundary of the normalized family.}
The missing Gaussian limit has a precise converse. In the weak topology
of log-spectral probability laws, the closure of the normalized mixtures
is exactly the class obtained by also allowing Gaussian kernels of every
scale. To see this, let
$\ell_\alpha=\frac12\log(D+2\alpha+1)$. For $G_D$ of gamma law
with shape $D$ and rate one, the gamma-ratio asymptotic in
\eqref{eq:ball-Mellin} gives
\[
 Y_\alpha-\ell_\alpha\ \Longrightarrow\ \tfrac12\log G_D
 \quad(\alpha\to\infty),\qquad
 (2\alpha+1)Y_\alpha\ \Longrightarrow\ \operatorname{Exp}(1)
 \quad(\alpha\downarrow-1/2).
\]
For the second formula the only nontrivial factor is
$\Gamma((\alpha+1/2)(1-it))/\Gamma(\alpha+1/2)\to(1-it)^{-1}$.
For $\alpha=\infty$ define $Y_\infty^{\circ}=\frac12\log G_D$,
and otherwise let $Y_\alpha^{\circ}=Y_\alpha-\ell_\alpha$.
The map $(\alpha,v)\mapsto\operatorname{Law}(Y_\alpha^{\circ}+v)$
on $(-1/2,\infty]\times\mathbb R$ is continuous and extends vaguely
by zero to its one-point compactification. This follows from uniform escape
at $-1/2$ proved above and uniform tightness of the centered profiles on
each $[a,\infty]$, $a>-1/2$. Its compact convex set of probability
barycenters on the compactification is therefore vaguely closed.
A limiting probability on the log line has full mass and consequently
places zero representing mass on the added zero generator. Every shape
at infinity gives a Gaussian up to scale, since its squared spectral
radius is a scaled $G_D$. Conversely, replacing every endpoint shape
by a finite shape tending to infinity while keeping $v$ fixed gives
weak convergence of any such mixture, by bounded convergence.
This proves the closure description.

The fixed-shape case $\alpha=0$ is Euclid's hat, characterized in all
dimensions by Gneiting, \emph{Radial positive definite functions
generated by Euclid's hat}, \emph{J. Multivariate Anal.} 69 (1999),
88--119, \href{https://doi.org/10.1006/jmva.1998.1800}{doi:10.1006/jmva.1998.1800}.
The case $d=1$, $\alpha=1$ has Gneiting's differential test, recalled
in \cite[Section 5.2]{JMR2009}. Mellin inversion for fixed-kernel
scale mixtures is classical. For a general inversion test under
Mellin convergence hypotheses see
J.~A. Palmer, \href{https://sccn.ucsd.edu/~jason/scalediff4.pdf}
{\emph{Differential criteria for positive definiteness}, Lemma~1}.
The compactness at both shape endpoints in \eqref{eq:ball-all}
is the additional step needed for the two-parameter family.

\section{A polynomial test for all scaled quartic Hermite extremals}
\label{sec:hermite-test}
Question 2(b) of Jaming, Matolcsi and R\'ev\'esz asks for a criterion
for membership in the scale mixtures of their extremal Hermite functions
of order four. Their Proposition 4.3 includes the Gaussian endpoint in
this family. The following test uses a single sequence of real numbers
computed from the candidate function. Each admissibility check on a test
polynomial requires only real-root isolation for univariate polynomials.

Let $f$ be continuous, even, real, bounded, $f(0)>0$, and $f(x)\to0$
as $|x|\to\infty$. Define a real linear functional on polynomials by
\[
 L_f(1)=1,\qquad L_f(X^k)=\frac1{(k-1)!f(0)}
 \int_0^\infty t^{k-1}e^{-t}f(\sqrt t)\,dt\quad(k\ge1).
\]
For $p\in\mathbb Q[X]$ write
\[
 B=X(1-X)p',\qquad
 C=(1-X)^2(X^2p''+2Xp').
\]
Call $p$ admissible if, for all $X\in[0,1]$ and $h\in[0,4/3]$,
\begin{align*}
 &p-2hB+h^2 C\ge0,\\
 &(1-h+3h^2/4)p+(2h-3h^2)B+h^2 C\ge0.
\end{align*}

\paragraph{Criterion.}
The function $f$ is a positive mixture over positive scales of the
quartic Hermite extreme ray generators if and only if $L_f(p)\ge0$
for every admissible rational polynomial $p$. Both rectangle tests are
effective. For $a(X)h^2+b(X)h+c(X)$ test its two endpoint values and,
where $a>0$ and $0<-b<8a/3$, also test $4ac-b^2\ge0$.
These sign conditions on a compact real interval are decided by isolating
the real roots of finitely many rational polynomials. Every failure of
membership has a certificate involving finitely many moments.
For a general real function, evenness, continuity, boundedness, a positive
value at zero and decay at infinity are necessary for any nonzero such
mixture and can be checked first. The zero function has the zero mixture.

\begin{proof}
In the variable $t=x^2$ the generators normalized at the origin, after
dilation, are $e^{-\lambda t}(1+b\lambda t+c\lambda^2t^2)$, $\lambda>0$,
where the two arcs in the $(b,c)$ plane are
\[
 (b,c)=(-2h,h^2),\qquad
 (b,c)=\left(\frac{2h-3h^2}{q_0},\frac{h^2}{q_0}\right),
 \quad q_0=1-h+3h^2/4,\quad 0\le h\le4/3 .
\]
To see this directly from the time and frequency alternatives in
Proposition 4.3 of the cited paper, a time-type extremal has polynomial
$(1-ht)^2$ because its four zeros in the spatial coordinate are real,
with even multiplicity and symmetry. Under the dimension-one Fourier
transform at the reference scale $t=\pi x^2$, this polynomial times its
Gaussian transforms to a Gaussian times
$q_0+(2h-3h^2)t+h^2t^2$. Here $q_0>0$. For $h\le2/3$ this polynomial
is nonnegative on $t\ge0$. For $h>2/3$ nonnegativity is equivalent
to $(2-3h)^2\le4q_0$, or $h\le4/3$. Conversely each time-type
polynomial with $h>0$ in this interval generates an extreme ray, since
all four of its zeros are real, as in Lemma 4.2 of the cited paper.
Fourier transformation gives the second arc, and $h=0$ gives the Gaussian.

Let $E$ be the union of these compact arcs, $Y=[0,1]\times E$, and
define $Tp(X,b,c)=p(X)+bB(X)+cC(X)$. For a kernel at scale $\lambda$,
with $X=(1+\lambda)^{-1}$, integration of the gamma moments gives for
$k\ge1$
\[
 L(X^k)=X^k\{1+kb(1-X)+k(k+1)c(1-X)^2\}=(TX^k)(X,b,c).
\]
This also holds for the constant polynomial. The origin values in the source's Hermite normalization form a continuous
strictly positive function on its compact parameter boundary, hence are bounded
above and away from zero. Normalization at the origin and denormalization
thus both preserve finiteness of the positive mixing measure.
The two admissibility tests
are precisely $Tp\ge0$ on $Y$. Mixtures satisfy all inequalities because
the normalized kernels and mixing measure are nonnegative.

Conversely, suppose that a normalized candidate passes. For any $m$ its
vector $(L_f(X^k))_{k=1}^m$ belongs to the compact convex hull of
$(TX^k(y))_{k=1}^m$, $y\in Y$. If it did not, strict separation,
approximation of the separating coefficients by rationals, and addition
of a rational constant would give $p\in\mathbb Q[X]$ with
$Tp>0$ and $L_f(p)<0$. Hence a probability measure on $Y$ matches these
first $m$ moments. The corresponding closed sets of probability measures
are nested in the weakly compact space of probability measures on $Y$,
so a single measure $\mu$ matches all moments.

For $0<X<1$ put $\lambda=(1-X)/X$, use the associated normalized
kernel in the variable $t>0$, and extend this kernel to the constant
$1$ at $X=1$ and the zero function at $X=0$. The normalized kernels
are bounded by $1$ by positive definiteness. Their $\mu$-average $v(t)$
is bounded and continuous on $t>0$, and it has every one of the
gamma moments of $u(t)=f(\sqrt t)/f(0)$. The Laplace transform of
$v-u$ is analytic on $\Re s>0$ and all of its derivatives vanish at
$s=1$. Analytic continuation and uniqueness yield $v=u$ on $t>0$.
As $t\downarrow0$ the limit of $v$ is $1-\mu(\{0\}\times E)$;
as $t\to\infty$ it is $\mu(\{1\}\times E)$. By continuity at zero
and the assumed decay, both endpoint atoms vanish. Thus $\mu$ represents
the required mixture over genuine positive finite scales. Conversely,
dominated convergence shows that both endpoint limits are necessary for
such a normalized mixture.
\end{proof}

\paragraph{A four-moment certificate.}
The degree-eight Hermite function $f(x)=(1-x^2)^4e^{-2x^2}$ is
nonnegative and positive definite, since it is the square of a quartic
time-type generator. Its exclusion from the mixture class is already
detected by
\[
 p(X)=10-160X+817X^2-1000X^3+433X^4,\qquad L_f(p)=-560/2187 .
\]
Indeed for $k\ge1$,
$L_f(X^k)=\sum_{j=0}^4(-1)^j\binom4j(k)_j3^{-k-j}$,
where $(k)_j$ is the rising factorial. The following finite certificate
verifies admissibility using rational arithmetic. Substitute $h=4Y/3$ in
the two tested polynomials $g(X,Y)$, of bidegree at most $(6,2)$.
On the $16$ by $16$ equal subdivision of the unit square, the respective
smallest bidegree-$(6,2)$ Bernstein coefficients over all boxes are
\[
 9413/49152\quad\hbox{and}\quad3693517/50331648.
\]
Here is an explicit recipe for verifying both numbers. For each
$i,j\in\{0,\ldots,15\}$ expand
$g((i+U)/16,(j+V)/16)=\sum_{k,l}d_{kl}U^kV^l$. For
$0\le a\le6$, $0\le b\le2$ take the minimum over all boxes and indices
of
\[
 \sum_{k=0}^a\sum_{l=0}^b d_{kl}
 \frac{\binom ak\binom bl}{\binom6k\binom2l}.
\]
Positivity of these coefficients implies positivity on the entire square
because the Bernstein basis functions are nonnegative and sum to one.

If a scale-mixture convention allows the constant kernel at
$\lambda=0$, simply drop the decay hypothesis in the criterion.
The same proof gives a finite nonnegative limit at infinity and identifies
it as the mass of that kernel times $f(0)$.

\subsection{Face dimensions at every Hermite order}
Write $\Omega_d^+$ for the cone of continuous radial functions on
$\mathbb R^d$ that are nonnegative and positive definite, and use the
Fourier kernel $e^{-2\pi i x\cdot\xi}$. Set $\delta=d/2$ and
$t=\pi|x|^2$. Define the polynomial involution $\mathcal H_d$ by
\[
 \widehat{e^{-t}P(t)}(\xi)=e^{-s}(\mathcal H_dP)(s),\qquad
 s=\pi|\xi|^2 .
\]
Its columns $G_k=\mathcal H_d(t^k)$ satisfy
$G_0=1$ and $G_{k+1}=(\delta+k-t)G_k+tG_k'$, by differentiating the
Fourier transform of $e^{-a\pi|x|^2}$ in $a$. It preserves polynomial
degree. A change of scale handles other Gaussian exponents.
Giraud and Peschanski, \emph{On positive functions with positive Fourier
transforms}, Acta Phys. Polon. B 37 (2006), 331--346, use Sturm sequences
to determine positivity in finite Hermite families, including the
self-Fourier family of degree eight. The face calculation below also uses
the polynomial domination established by Jaming, Matolcsi and R\'ev\'esz
in their Proposition 4.1 from Hardy's uncertainty principle.

\paragraph{A rank formula.}
Let $0\ne f=e^{-t}P(t)\in\Omega_d^+$ and $N=\deg P$. For every
nonnegative real zero $a$ of $P$, of multiplicity $m$, take the rows
$Q\mapsto Q^{(j)}(a)$ for $0\le j<m$. For each such zero $a$ of
$\mathcal H_dP$ take the rows $Q\mapsto(\mathcal H_dQ)^{(j)}(a)$
with the corresponding multiplicity. On $\mathbb R[t]_{\le N}$, let
$M$ be the matrix of all these rows. The minimal face of $\Omega_d^+$
containing $f$ has dimension $N+1-\operatorname{rank}M$.
In particular, $f$ is extreme exactly when $\operatorname{rank}M=N$.

\begin{proof}
If $g$ and $Cf-g$ belong to $\Omega_d^+$ for $C>0$, the functions
$g$ and $\widehat g$ are bounded above by the corresponding Gaussian
times a polynomial of spatial degree at most $2N$. Here integrability
of $g$ follows by domination and nonnegativity of both Fourier transforms
follows by positive definiteness. To reduce the radial case to the
classical one-dimensional polynomial version of Hardy's theorem, project
$g$ onto the first coordinate by integrating over its orthogonal
complement. This projection has a bound
$C'(1+|u|)^{2N}e^{-\pi u^2}$, and its one-dimensional Fourier transform
equals $\widehat g(\xi e_1)$ with the same type of bound. Thus the
restriction $\widehat g(\xi e_1)$ is a Gaussian times an even polynomial
of degree at most $2N$. Radiality and Fourier inversion now give
$g=e^{-t}Q(t)$ with $\deg Q\le N$. On both sides $Q$ and
$\mathcal H_dQ$ vanish at the specified zeros to at least the indicated
multiplicities, since they are bounded between zero and the corresponding
multiples of the two nonnegative polynomials. Hence $MQ=0$.

Conversely, if $MQ=0$, the ratios $Q/P$ and
$(\mathcal H_dQ)/(\mathcal H_dP)$ are bounded on $[0,\infty)$.
Their real singularities are removable and their degrees bound them
at infinity. Therefore $f\pm\epsilon e^{-t}Q(t)$ belong to $\Omega_d^+$
for sufficiently small $\epsilon>0$. The elementary description of the
minimal face of a cone at $f$ as the set of nonnegative cone elements
dominated by some positive multiple of $f$ shows that its linear span
is exactly $e^{-t}\ker M$. This proves the dimension formula.
\end{proof}

\paragraph{A degree-eight mixed example in every dimension.}
In each of the following cases let $P(t)=(t-a)^2Q(t)$.
\[
\begin{array}{c|c|l}
d&a&Q(t)\\ \hline
1,2,3 &(\delta+1)/3&
 t^2-4(\delta+4)t/3+(4\delta^2+29\delta+79)/9\\
d\ge4 &\delta/2&
 t^2-(\delta+6)t+\delta^2/4+5\delta+6
\end{array}
\]
Then $e^{-t}P(t)$ is an extreme ray generator of $\Omega_d^+$.
In both rows the recurrence for $G_k$ gives $\mathcal H_dP=P$.
The discriminants of $Q$ in the two rows are $4(\delta-5)/3$ and
$12-8\delta$, respectively, and are negative on the indicated ranges.
There is exactly one nonnegative zero $a$ with multiplicity two, on
each side. The four rows of $M$ are evaluation and first derivative
at $a$ in time and frequency. Applying them to $1,t,t^2,t^3$ and
using the same recurrence gives the respective determinants
\[
 -\frac{(\delta+1)^2(\delta^2+2\delta+28)}{81},
 \qquad -2\delta .
\]
Both are nonzero, giving rank four. The two time conditions alone, and
likewise the two frequency conditions alone, have rank two. For $d=1$ a
scalar multiple of the displayed polynomial is
$(2t-1)^2(2t^2-12t+21)$. Thus the all-real-zeros sufficient condition
of Jaming, Matolcsi and R\'ev\'esz does not certify this example.

\paragraph{Extremality against nonradial summands.}
The same examples are extreme in the cone of all continuous nonnegative
positive definite functions on $\mathbb R^d$. For $d=1$ real-valued
positive definite functions are even. Assume $d\ge2$. For a summand $g$,
apply the projection argument in every unit direction. The entire function
$e^{\pi\zeta\cdot\zeta}\widehat g(\zeta)$ restricts on each real line
through zero to a polynomial of degree at most eight. All its homogeneous
Taylor coefficients of higher degree vanish, by restriction to real unit
vectors. Fourier inversion therefore gives $g=e^{-t}D(x)$ with $D$ a
real even polynomial of degree at most eight.
Expand $D$ in homogeneous harmonic polynomials $Y_{\ell j}(x)$ times
radial polynomials $q_{\ell j}(t)$ of degree at most $(8-\ell)/2$.
Only even $\ell$ occur. Under Fourier transformation each term keeps
its harmonic polynomial and its radial factor becomes
$(-i)^\ell e^{-s}\mathcal H_{d+2\ell}q_{\ell j}(s)$.
This follows by differentiating in the scale parameter the identity
\[
 \widehat{Y_{\ell j}(x)e^{-b\pi|x|^2}}(\xi)
 =(-i)^\ell b^{-d/2-\ell}Y_{\ell j}(\xi)e^{-\pi|\xi|^2/b};
\]
the identity follows by applying the harmonic differential operator
$Y_{\ell j}(\partial_\xi)$ to the Gaussian transform.

Domination by $f$ and $\widehat f$ forces a double radial zero at $a$ on
each ray on both sides, so spherical harmonic orthogonality imposes the
same two conditions on each $q_{\ell j}$ and its transform.
Degrees $\ell\ge6$ leave radial polynomials of degree at most one and
these vanish. For $\ell=4$ the time constraints leave only multiples of
$(t-a)^2$; the value of $\mathcal H_{d+8}((t-a)^2)$ at $a$ is
$(\delta+10)(\delta+13)/9$ in the first row and $20$ in the second.
For $\ell=2$ the four constraints on $1,t,t^2,t^3$ have determinant
\[
 -\frac{(\delta+7)(\delta^3+21\delta^2+174\delta+316)}{81}
 \quad\hbox{or}\quad -6(\delta+12),
\]
respectively, by the recurrence with $\delta$ replaced by $\delta+2$.
Thus all nonconstant angular components vanish and the radial calculation
forces $g$ to be proportional to $f$.

More generally the same argument computes the dimension of the minimal
face of the full, unrestricted cone at any radial Hermite function above.
For $d\ge2$ and even $0\le\ell\le2N$ set $n_\ell=N-\ell/2$.
Form $M_\ell$ on $\mathbb R[t]_{\le n_\ell}$ with the same physical zeros
of $P$ and frequency zeros of $\mathcal H_dP$ as before, replacing
$\mathcal H_dQ$ in the frequency rows by $\mathcal H_{d+2\ell}Q$.
If $h_{d,\ell}$ denotes the real dimension of the spherical harmonics
of degree $\ell$, the dimension is
\[
 \sum_{\substack{0\le\ell\le2N\\\ell\ {\rm even}}}
 h_{d,\ell}\big(n_\ell+1-\operatorname{rank}M_\ell\big).
\]
Indeed every summand obeys all sector constraints by the projection and
harmonic decomposition just used. Conversely any element of these kernels
produces a perturbation bounded by a multiple of $f$ on both sides:
compact angular spheres bound each harmonic factor, vanishing orders
control finite radii and the total degree controls infinity. The origin
creates no exceptional constraint because $P(0)=f(0)>0$ and
$(\mathcal H_dP)(0)=\int f>0$. Small two-sided perturbations therefore
give exactly the asserted linear span.

There are also two-parameter families of such rays without Fourier
self-duality. Near each displayed pair $(a,a)$, prescribe a time double
zero at $u$ and a frequency double zero at $v$. The four time and frequency
conditions determine a unique monic quartic $P_{u,v}$ in $t$, since
the minor on $1,t,t^2,t^3$ is nonzero at $(a,a)$ and hence nearby.
The coefficients depend continuously on $(u,v)$. Both residual monic
quadratics, obtained by dividing $P_{u,v}$ by $(t-u)^2$ and
$\mathcal H_dP_{u,v}$ by $(t-v)^2$, keep negative discriminants.
Thus both sides are nonnegative. The nonzero angular-sector minors just
computed remain nonzero in this neighborhood (for $\ell=4$ use the
nonzero value of the transformed time quadratic at $v$).
Consequently all these rays are extreme in the full cone. When $u\ne v$
the zeros on the two sides differ, so the function is not self-Fourier.
For an explicit one-dimensional pair take $u=1/2$, $v=1$ and
\[
 P=(2t-1)^2(4t^2-36t+95),\qquad
 \mathcal H_1P=8(t-1)^2(2t^2-4t+11).
\]
The two quadratic discriminants are negative and the minor on
$1,t,t^2,t^3$ is $-3/4$, proving extremality directly.

\section{A criterion that also detects signed functions}
\label{sec:signed}
Question 2(c) of \cite{JMR2009} asks for a positive-definiteness criterion outside mixtures of nonnegative positive definite functions. Here is a one-time Gaussian moment test. It applies to continuous bounded real radial $f$ on $\mathbb R^d$ without an integrability or a sign assumption. It uses classical moment theory; its inclusion specifies how this part of the question can be met within the same Fourier framework.

Put $D=d/2$, $v=|x|^2/4$, and compute the real numbers
\[
 m_n=\frac{n!}{(4\pi)^D}\int_{\mathbb R^d}e^{-v}L_n^{D-1}(v) f(x)\,dx\qquad(n\ge0),
\]
where $L_n^{D-1}$ is the generalized Laguerre polynomial. Then $f$ is positive definite if and only if, for every $N\ge0$, both Hankel matrices $(m_{j+k})_{j,k=0}^N$ and $(m_{j+k+1})_{j,k=0}^N$ are positive semidefinite. For instance this tests the signed function $\cos x$ on the line, whose moments here all equal $e^{-1}$.

To prove the criterion, let $G(s)=(4\pi s)^{-D}\int e^{-|x|^2/(4s)}f(x)\,dx$. This is holomorphic for $\Re s>0$ by locally uniform Gaussian domination. The Laguerre generating function, or differentiation under the integral, gives $(-1)^nG^{(n)}(1)=m_n$. For a positive definite $f$, let $\mu$ be the pushforward of its Bochner measure under $\xi\mapsto |\xi|^2$, with phase $e^{ix\cdot\xi}$. Then $G(s)=\int e^{-su}\,d\mu(u)$ and $m_n$ are the Stieltjes moments of $e^{-u}\mu$, giving the two matrix inequalities.

Conversely the Stieltjes moment theorem supplies a nonnegative measure $\eta$ on $[0,\infty)$ with these moments \cite{Akhiezer}. For every $r<1$, $r>0$, analyticity gives $m_n\le C_r n!r^{-n}$. Monotone summation of the exponential series then gives $\int e^{vu}\,d\eta(u)<\infty$ for $0\le v<1$. Set $d\mu=e^u d\eta$. Comparison of Taylor series at 1 and analytic continuation give $G(s)=\int e^{-su}\,d\mu(u)$ for $s>0$. As $s\downarrow0$, the Gaussian approximate identity gives $G(s)\to f(0)$, so $\mu$ is finite. Distribute its mass at each $u$ uniformly on the sphere of radius $\sqrt u$. The Fourier transform of the resulting measure is a continuous radial positive definite function $g$ with the same $G(s)$. Uniqueness of the ordinary Laplace transform in $v=|x|^2/4$, with radial weight $v^{D-1}$, forces $f=g$ first almost everywhere and then everywhere by continuity.

\section{The optimal number of half-ball autocorrelations}
\label{sec:factor-rank}
Let $B_R$ be the closed Euclidean ball of radius $R>0$ centered at the origin of $\mathbb R^d$, and put $\widetilde u(x)=\overline{u(-x)}$. Given a continuous radial positive definite function $f$ supported in $B_R$, consider representations
\begin{equation}\label{eq:factor-rank}
 f=\sum_{j=1}^{m}u_j*\widetilde u_j,
 \qquad u_j\in L^2(\mathbb R^d),\qquad \operatorname{supp}u_j\subseteq B_{R/2}.
\end{equation}
The factors may be nonradial. Ehm, Gneiting and Richards characterize the existence of a single factor in terms of the zeros of a univariate entire function \cite[Theorems~2.1 and 3.1]{EGR}. They state a countable real representation and also give $d$ terms for a class of extreme points \cite[Theorems~3.8--3.9 and the ensuing discussion]{EGR}. The countable radial scalar and gradient expansion is also stated in \cite[eqs.~(39)--(40)]{EGR}. Rudin's smooth expansion and Efimov's continuous extension in dimensions $d\ge3$ are discussed in \cite[\S1 and Remark~7]{GM}. We ask for the smallest number of terms that suffices uniformly in $f$.

\begin{theorem}\label{thm:sharp-factor-number}
For every $d\ge1$ the least universal value of $m$ in \eqref{eq:factor-rank} is $d$ for real-valued factors and $\lfloor d/2\rfloor+1$ for complex-valued factors. In each dimension a single covariance attains both lower bounds.
\end{theorem}
\begin{proof}
Use the Fourier transform with phase $e^{-ix\cdot\xi}$. The spectral density $\widehat f(\xi)=q(|\xi|)$ is nonnegative and integrable by Bochner's theorem and Fourier uniqueness, since $f$ is integrable and continuous at zero. The function $q(s)=\widehat f(se_1)$ is even and entire, with $|q(s)|\le C e^{R|\operatorname{Im}s|}$. It belongs to $L^1(\mathbb R)$ by polar integration at infinity and boundedness near zero. In fact its inverse one-dimensional Fourier transform is $t\mapsto\int_{\mathbb R^{d-1}} f(t,y)\,dy$, a real even continuous positive definite function supported in $[-R,R]$. The one-dimensional real-factor theorem \cite[Theorem~2.1(b)]{EGR} therefore gives
\[
 q(s)=|h(s)|^2\quad(s\in\mathbb R),\qquad
 |h(z)|\le Ce^{R|\operatorname{Im}z|/2},\qquad h(-s)=\overline{h(s)}\quad(s\in\mathbb R),
\]
where $h$ is the Fourier transform of a real $L^2$ function supported in $[-R/2,R/2]$. Write $h(s)=E(s)+isP(s)$, so $E(s)=(h(s)+h(-s))/2$ and $P(s)=(h(s)-h(-s))/(2is)$ with the removable value at zero. Both $E,P$ are real on the real axis, even and entire. Both obey the exponential bound for $h$, by division outside the unit disk and boundedness inside it.

For $z\in\mathbb C^d$ take $w^2=\sum_{j=1}^d z_j^2$. The even power series define entire functions $E(w)$ and $P(w)$ independently of the choice of square root. If $z=x+iy$, then $|\operatorname{Im}w|^2=(|\sum_j z_j^2|-|x|^2+|y|^2)/2\le|y|^2$. Thus $E(w)$ and $z_jP(w)$ have bounds $C(1+|z|)e^{R|\operatorname{Im}z|/2}$. On the real axes, with $r=|\xi|$,
\[
 q(r)=E(r)^2+\sum_{j=1}^{d}\xi_j^2P(r)^2 .
\]
All the displayed summands are squares of $L^2(\mathbb R^d)$ functions by spectral integrability. Paley--Wiener--Schwartz and Plancherel give half-ball support for the inverse transform of each of these square roots and for their linear combinations.

For real spatial factors take the $d$ Fourier transforms $E(r)+i\xi_dP(r)$ and $i\xi_jP(r)$ for $1\le j<d$. They are Hermitian and their squared moduli sum to $q(r)$. For complex factors, when $d$ is even, take $E(r)$ and the $d/2$ functions $(\xi_{2j-1}+i\xi_{2j})P(r)$. When $d$ is odd, use $E(r)+i\xi_dP(r)$ and the same pairs for the other $d-1$ coordinates. Taking inverse transforms proves both upper bounds.

For sharpness choose an integer $M>(d+2)/2$, set $\alpha=R/(2M)$, and put $S(t)=(\sin(\alpha t)/(\alpha t))^M$, with $S(0)=1$. Consider
\[
 q_0(r)=(1+r^2)S(r)^2 .
\]
Polynomial times even powers of sinc also occur in the examples of \cite[eq.~(18)]{EGR}. The profiles $S(w)$ and $z_jS(w)$ have the preceding half-ball Paley--Wiener bounds and square-integrable real restrictions, by the choice of $M$. Their squared moduli sum to $q_0$ on $\mathbb R^d$. Thus $q_0$ is the spectral density of a continuous radial positive definite function supported in $B_R$, and it may be normalized by multiplying by a positive constant.

Suppose $q_0(|\xi|)=\sum_{j=1}^m|H_j(\xi)|^2$ for Fourier transforms of arbitrary half-ball factors. Each is entire with $|H_j(z)|\le C_j e^{R|\operatorname{Im}z|/2}$ by Cauchy--Schwarz on the ball. Fix a real unit vector $\theta$. On the real axis $|H_j(t\theta)|\le (1+t^2)^{1/2}|S(t)|$, so $H_j(t\theta)$ vanishes to order at least $M$ at every nonzero $t\in(\pi/\alpha)\mathbb Z$. The quotient $Q_j(z)=H_j(z\theta)/S(z)$ is therefore entire. Outside fixed disjoint disks around the points of $(\pi/\alpha)\mathbb Z$ one has $|\sin(\alpha z)|\ge c e^{\alpha|\operatorname{Im}z|}$, by periodicity on a bounded horizontal strip and the exponential formula off that strip. The choice $M\alpha=R/2$ now gives $|Q_j(z)|\le C(1+|z|)^M$ there. The maximum principle fills the disks with the same polynomial bound up to a constant. By Cauchy's estimates $Q_j$ is a polynomial, and $|Q_j(t)|\le(1+t^2)^{1/2}$ forces it to be affine. Since $S'(0)=0$, its coefficients are $H_j(0)$ and $\nabla H_j(0)\cdot\theta$. Varying $\theta$ shows that, for some vectors $c,b_1,\ldots,b_d\in\mathbb C^m$,
\[
 (H_1(\xi),\ldots,H_m(\xi))=S(|\xi|)\biggl(c+\sum_{j=1}^d\xi_j b_j\biggr).
\]
Comparing the squared norm with $q_0$ shows that $c,b_1,\ldots,b_d$ are orthonormal in the underlying real Euclidean space $\mathbb R^{2m}$. Hence $2m\ge d+1$. For real spatial factors, Hermitian symmetry additionally puts all $b_j$ in $i\mathbb R^m$. Their orthonormality forces $m\ge d$. This proves both sharp bounds.
\end{proof}

\section{From a half-ball factorization to spherical means}
\label{sec:sered}
Let $B$ denote the closed Euclidean ball of radius $1/2$. For $d\ge2$ and $0<a<1$ let $\sigma_a$ be the probability measure on the sphere of radius $a$. Efimov's objective $\Phi_d(a)$ is the supremum of $|\int f\,d\sigma_a|$ over continuous real $f$ supported in $2B$ with $f(0)=1$ and $\widehat f\ge0$ \cite{EfimovSphere,GM}. The radial spectral reduction is already known. We recall it briefly to connect the sharp factor count with the calculations that follow.

Let $A$ on $L^2(B)$ denote convolution with $\sigma_a$ after extension by zero, followed by restriction to $B$. It is self-adjoint by symmetry. A unit factor $u$ gives an admissible autocorrelation for a real $u$ and taking the real part works for complex $u$; its spherical integral is the Rayleigh quotient $\langle Au,u\rangle$. Conversely rotational averaging of any admissible $f$ preserves its normalization and objective. By Theorem~\ref{thm:sharp-factor-number} this radial average is a sum of half-ball autocorrelations, with $\sum_j\|u_j\|_2^2=f(0)=1$. Consequently $\Phi_d(a)=\|A\|$. There are nonnegative radial near-maximizers on the factor side: replace $u(x)$ by $(\int_{O(d)}|u(Ux)|^2\,dU)^{1/2}$ using normalized Haar measure. Rotation invariance and Cauchy--Schwarz inside the rotation integral bound the absolute Rayleigh numerator by the new one and preserve the norm. The spectral construction below gives attainment.

\section{The three-dimensional spherical Tur\'an problem}
\label{sec:efimov}
Let $B$ be the unit Euclidean ball in $\mathbb R^3$, and let $\sigma_a$ be probability surface measure on the sphere of radius $a$. Efimov's original constant is
\[
 \Phi_3(a)=\sup\left\{\left|\int f\,d\sigma_a\right|:\ f\in C(\mathbb R^3,\mathbb R),\quad
 f(0)=1,\quad \operatorname{supp}f\subseteq B,\quad \widehat f\ge0\right\},\qquad 0<a<1.
\]
Efimov reduced the problem to an integral operator and evaluated it for $1/3\le a<1$ \cite[Theorems 3--5]{EfimovSphere}. We obtain a finite scalar equation for the full interval, including the range $0<a<1/3$. Its matrix order is $\lfloor 1/a\rfloor$. It includes the reciprocal-integer transition radii without separate limits. Related finite systems for an inhomogeneous moving-average equation occur in the work of Freeman and Miller on one-level densities \cite{Freeman}.

For $m\ge1$ index matrix coordinates from $0$ to $m-1$. Let $J_m$ have $+1$ on the superdiagonal, $-1$ on the subdiagonal and zero elsewhere. Let $E_m$ have columns $e_j-e_{m-1-j}$ for $0\le j<\lfloor m/2\rfloor$. It has no columns when $m=1$. Set $n=\lfloor 1/a\rfloor$, $d=1-na$, and let $P:\mathbb R^{n+1}\to\mathbb R^n$ discard the last coordinate. Define
\begin{equation}
 D_a(z)=\det\!\left[\ P e^{zdJ_{n+1}/2} E_{n+1}\ \big|\ -e^{-z(a-d)J_n/2} E_n\ \right].
 \label{eq:ef-det}
\end{equation}
The column counts sum to $n$. The eigenvalues of $J_m$ are $2i\cos(k\pi/(m+1))$, $k=1,\ldots,m$, as follows by diagonal unitary conjugation from the adjacency matrix of a path. Thus \eqref{eq:ef-det} uses only finite-dimensional linear algebra and trigonometric functions of specified arguments.

\begin{theorem}\label{thm:ef-all}
For every $a\in(0,1)$, $D_a(0)\ne0$, a positive zero exists, and if $z_1(a)$ is its smallest positive zero then
\begin{equation}
  \Phi_3(a)=\frac1{2a z_1(a)}.
  \label{eq:ef-value}
\end{equation}
All extremal values in this formula are attained. A null vector of the matrix in \eqref{eq:ef-det} at $z=z_1(a)$ constructs an attaining normalized autocorrelation as described below.
\end{theorem}

\begin{proof}
Efimov proved that an extremizer has the form $f=u*\widetilde u$, where $u$ is a real radial $L^2$ function supported in the radius-$1/2$ ball and $\|u\|_2=1$. Conversely, every such autocorrelation is admissible. Write $L=1/2$, $v(r)=r u(r)$ and extend $v$ oddly across the origin and by zero outside $[-L,L]$. The map to $v$ is an isometry up to a fixed constant. The average of $u(x+a\omega)$ over the unit sphere, at $|x|=r>0$, is
\[
 \frac1{2ar}\int_{|r-a|}^{r+a}s u(s)\,ds.
\]
Indeed, the coordinate of $\omega$ along $x$ is uniform on $[-1,1]$. Oddness converts the integral into $\int_{r-a}^{r+a}v(s)\,ds$. Consequently the extremal problem is the absolute Rayleigh maximum for $(2a)^{-1}K_a$ on odd $v$, where $K_a v(r)=\int_{r-a}^{r+a}v(s)\,ds$ on $[-L,L]$. This is a compact self-adjoint operator, as its kernel on the bounded interval is square integrable. On the positive half-interval the odd kernel is
\[
 \mathbf1_{\{|r-s|<a\}}-\mathbf1_{\{r+s<a\}}
   =\mathbf1_{\{|r-s|<a<r+s\}}\quad\text{almost everywhere}.
\]
It is nonnegative and positive on a nonempty open patch since $a<2L$. The absolute value of a Rayleigh quotient is bounded by the Rayleigh quotient for the pointwise absolute value on the half-interval. Hence the norm is a positive eigenvalue, equal to the largest positive eigenvalue. Neither a sign hypothesis on the original $f$ nor radiality of its admissible class has been imposed.

Translate to $[0,1]$ by $w(x)=v(x-1/2)$, so that $w(1-x)=-w(x)$. Extend by zero outside the interval. For $z\ne0$ the equation $w=zK_aw$ implies the delay equation
\begin{equation}
 w'(x)=z\,[w(x+a)-w(x-a)]\qquad\text{for a.e. }x\in(0,1).
 \label{eq:ef-delay}
\end{equation}
Indeed the moving average is absolutely continuous. Conversely, if $w$ is absolutely continuous inside the closed interval, odd about $1/2$, and satisfies \eqref{eq:ef-delay}, then $w-zK_aw$ is constant. Its oddness makes that constant zero. No condition on the one-sided endpoint values is required. The zero extension need not be continuous at the two outer endpoints.

For $0<t<d$ write $p(t)=(w(ja+t))_{j=0}^{n}$ and for $d<t<a$ write $q(t)=(w(ja+t))_{j=0}^{n-1}$. The two systems are $p'=zJ_{n+1}p$ and $q'=zJ_nq$. Reversal of coordinates anticommutes with $J_m$. Reflection about the midpoint of $[0,1]$ consequently means that the vectors at the two strip midpoints belong to the ranges of $E_{n+1},E_n$. Explicitly,
\[
 p(t)=e^{z(t-d/2)J_{n+1}}E_{n+1}u_0,\qquad
 q(t)=e^{z(t-(a+d)/2)J_n}E_n u_1 .
\]
These formulae satisfy $p_j(t)=-p_{n-j}(d-t)$ and $q_j(t)=-q_{n-1-j}(a+d-t)$. The $n$ conditions $Pp(d)=q(d)$ are exactly the matrix equation in \eqref{eq:ef-det}. Reflection then gives $q_j(a)=p_{j+1}(0)$ for $j<n$. Thus the pieces match everywhere in the open interval. When $d=0$, $p(0)$ denotes the $n+1$ knot values only. The same matrix matches $q(0)$ to its first $n$ coordinates, and reflection matches $q(a)$ to its last $n$ coordinates. A nonzero null vector always yields a nonzero function. In the case $d=0$ this uses injectivity of $P$ on the antisymmetric subspace, and otherwise it follows from the two invertible propagators on positive-length strips.

This constructs an equivalence between nonzero null vectors and eigenfunctions for eigenvalue $1/z$ of $K_a$ on the odd subspace. For $z=0$, any matching solution would be a continuous odd constant and hence zero, proving $D_a(0)\ne0$. Since the norm eigenvalue is positive, there is a positive zero. By compactness of $K_a$ its largest positive eigenvalue corresponds precisely to the smallest positive zero. For that zero, reverse the shift, put $u(r)=v(r)/r$ on $(0,L)$, extend radially by zero, and normalize in $L^2(\mathbb R^3)$. This completes the construction and the proof.
\end{proof}

For $a=1/4$, put $t=z/8$ and let $t_*$ be the unique solution in $(0,\pi/(2\sqrt5))$ of
\begin{equation}
  \frac3{\sqrt5}\sin(\sqrt5\,t)=\cos t.
  \label{eq:ef-quarter}
\end{equation}
Then $\Phi_3(1/4)=1/(4t_*)$, approximately $0.719832$. To verify the equation, here $n=4$ and $d=0$. The nonconstant part of the determinant in \eqref{eq:ef-det} uses two columns of $e^{-tJ_4}$, hence lies in the span of $1,\cos t,\sin t,\cos(\sqrt5t),\sin(\sqrt5t)$. Indeed the eigenvalues on the second exterior power of the exponential are the exponentials of pairwise sums of the four eigenvalues of $-tJ_4$. Those of $J_4$ are $\pm i(1+\sqrt5)/2$ and $\pm i(\sqrt5-1)/2$. Direct expansion of the four-by-four matrix to order four gives the derivatives of $D_{1/4}(8t)$ at zero as $-1,3,1,-15,-1$, respectively. Thus it equals $-\cos t+(3/\sqrt5)\sin(\sqrt5t)$. The left side of \eqref{eq:ef-quarter} increases strictly and the right side decreases strictly on the displayed interval. The signs at the endpoints give the first positive root as asserted.

\section{Finite equations in every odd dimension}
\label{sec:efimov-odd}
The finite reduction extends beyond the box kernel. In dimension $d=2k+1$, $k\ge1$, the projection of uniform surface measure onto a line has polynomial density on its support. Passing from a three-dimensional odd function to a one-dimensional function with prescribed parity and moments then gives a finite characteristic equation for the original spherical objective $\Phi_{2k+1}(a)$ at every radius $0<a<1$. The scalar reduction below agrees with the general spectral framework of Gorbachev and Manov \cite[Sec.~5, eq.~(5.5)]{GM}. We specify all entries of the finite system.

Let $I=[-1/2,1/2]$ and use $\widehat h(s)=\int h(x)e^{-isx}dx$. Put $\epsilon=(-1)^k$ and let $S_k$ be the closed subspace of real $L^2(I)$ defined by
\[
 v(-x)=\epsilon v(x),\qquad \int_I x^jv(x)\,dx=0\quad (0\le j<k).
\]
The map taking a real radial $u\in L^2(\mathbb R^{2k+1})$ supported in the half-ball to the one-dimensional $v$ defined by
\[
 \widehat v(s)=(is)^k\widehat u(se_1)
\]
is a bijection onto $S_k$ and preserves the norm up to a constant depending only on $k$. Indeed, $F(s)=\widehat u(se_1)$ is even, entire and bounded by $C e^{|\operatorname{Im}s|/2}$, since $u$ is also integrable. Radial Plancherel shows $s^k F(s)\in L^2(\mathbb R)$. One-dimensional Paley--Wiener gives the required support of $v$, and the zero at $s=0$ gives its moment constraints. Conversely, for $v\in S_k$, $F(s)=\widehat v(s)/(is)^k$ is even and entire, real on $\mathbb R$, and satisfies the same exponential bound by removable continuation at zero. The even power series defines $F(\sqrt{z\cdot z})$ on $\mathbb C^{2k+1}$, with exponential bound $C e^{|\operatorname{Im}z|/2}$ because $|\operatorname{Im}\sqrt{z\cdot z}|\le|\operatorname{Im}z|$. Its restriction to $\mathbb R^{2k+1}$ belongs to $L^2$ by polar integration. Paley--Wiener and Plancherel give back the required real radial half-ball factor.

The projected spherical density is $p(t)\mathbf1_{(-a,a)}(t)$, where
\[
 p(t)=\frac{\Gamma(k+1/2)}{a\sqrt\pi\,\Gamma(k)}(1-t^2/a^2)^{k-1}.
\]
Write $C_a v=(p\mathbf1_{(-a,a)})*v$ restricted to $I$, with $v$ zero outside $I$. Efimov's attainment by a real radial autocorrelation \cite[Theorem~3]{EfimovSphere}, radial Plancherel and the preceding isometry give
\begin{equation}
 \Phi_{2k+1}(a)=\left\|\left.\Pi_{S_k} C_a\right|_{S_k}\right\|,
 \label{eq:odd-compression}
\end{equation}
where $\Pi_{S_k}$ is orthogonal projection. In detail, the radial Fourier multiplier of spherical averaging is $\widehat{p\mathbf1_{(-a,a)}}(|\xi|)$, so multiplication on the frequency side gives the Rayleigh quotient of $C_a$ for the corresponding $v$. This compression is compact self-adjoint since the restricted kernel is square integrable. Its norm is strictly positive, as the normalized half-ball autocorrelation is strictly positive on every sphere of radius $a<1$. Using the eigenvalue of greatest absolute value below avoids any requirement that the compressed kernel have a fixed sign.

We give the explicit finite system whose loss of rank yields that norm. Shift the coordinate on $I$ to $x\in[0,1]$ and continue to write $v(x)$ for the shifted function. Let $N=2k-1$, $n=\lfloor 1/a\rfloor$ and $d_0=1-na$. Use the following nonempty strips, omitting the first when $d_0=0$:
\[
 (l,h,m)=(0,d_0,n+1),\qquad (d_0,a,n).
\]
Each such strip carries $y_j(t)=v(ja+t)$ for $l<t<h$ and $0\le j<m$. The cells $ja+(l,h)$ partition $(0,1)$ up to their endpoints. For a real parameter $z$ the permitted functions on a strip are the solutions of the constant-coefficient system
\begin{equation}
 y_j^{(N)}(t)=z\sum_{s=0}^{N-1}\left[p^{(N-1-s)}(-a)y_{j+1}^{(s)}(t)
       -p^{(N-1-s)}(a)y_{j-1}^{(s)}(t)\right],
 \label{eq:odd-ode}
\end{equation}
where the two missing neighbours are zero. Regard the $mN$ midpoint values of $y_j^{(s)}$, $0\le j<m$, $0\le s<N$, as independent initial coordinates. For fixed $z$ all $y_j(t)$ are explicit entries of a matrix exponential times this initial vector. Add further independent unknowns for a polynomial
\[
 Q(x)=\sum_{\substack{0\le j<k\\(-1)^j=\epsilon}} c_j(x-1/2)^j.
\]
On this finite vector of unknowns impose the following homogeneous real linear equations.
\begin{enumerate}
\item On each strip and for each $j,s$, set the midpoint value of $y_j^{(s)}$ equal to $\epsilon(-1)^s$ times that of $y_{m-1-j}^{(s)}$.
\item Require $\int_0^1(x-1/2)^jv(x)\,dx=0$ for $0\le j<k$.
\item On each cell $ja+(l,h)$ evaluate
\[
       v(x)-z(C_a v)(x)=Q(x)
\]
at the $N$ distinct nodes $x=ja+l+i(h-l)/(N+1)$, $1\le i\le N$.
\end{enumerate}
Call the resulting rectangular matrix $S_{k,a}(z)$ and define
\[
 \mathcal D_{k,a}(z)=\det\bigl(S_{k,a}(z)^{\mathsf T}S_{k,a}(z)\bigr).
\]
Every entry is a finite expression. For each fixed node in the last condition, split $[x-a,x+a]\cap[0,1]$ along the cell endpoints. The integrand on each piece is a polynomial times a matrix exponential. For example all integrals $\int_0^t s^j e^{sA}ds$ are obtained from a single block-triangular matrix exponential by augmenting the constant coefficient system for $e^{sA}$ with those for $s^i e^{sA}$ and their primitives, $0\le i\le j$. This prescription remains valid when the coefficient matrix is singular.

\begin{theorem}\label{thm:odd-finite}
For every $k\ge1$ and $0<a<1$, $\mathcal D_{k,a}(0)>0$. Its nonzero real roots are precisely the reciprocals of the nonzero eigenvalues in \eqref{eq:odd-compression}. Consequently
\[
 \Phi_{2k+1}(a)=\left(\min\{|z|:z\in\mathbb R,\ \mathcal D_{k,a}(z)=0\}\right)^{-1}.
\]
A null vector at any minimizing root yields an extremizer by reconstructing $v$ from the strip solutions, applying the inverse radial transform above, and taking the normalized autocorrelation.
\end{theorem}
\begin{proof}
Within its parity subspace, $S_k$ is the orthogonal complement of the parity-$\epsilon$ polynomials of degrees below $k$. The eigenvalue equation for $1/z$, $z\ne0$, therefore says exactly that $v\in S_k$ and $v-zC_av=Q$ for some such polynomial. In the sense of distributions, since $p$ has degree $N-1$,
\[
 (p\mathbf1_{(-a,a)})^{(N)}
 =\sum_{s=0}^{N-1}\bigl[p^{(N-1-s)}(-a)\delta_{-a}^{(s)}
                    -p^{(N-1-s)}(a)\delta_{a}^{(s)}\bigr].
\]
Taking $N$ derivatives on cell interiors thus gives \eqref{eq:odd-ode}; $Q^{(N)}=0$. The leading coefficient matrix of this system is the identity. Its distributional solutions on each open strip are classical entire exponential-polynomial solutions there, as seen by multiplying the corresponding first-order system by the inverse fundamental matrix. Reflection preserves the system, since its scalar convolution equation has an even kernel (equivalently, use the parity of the derivatives of $p$ at the two endpoints). Thus midpoint conditions enforce the stated parity throughout each strip. Every actual eigenfunction satisfies all three sets of finite conditions.

Conversely, the strip ODE makes the $N$th derivative of $v-zC_av-Q$ vanish on each cell in the distributional sense. The residual therefore is a polynomial of degree less than $N$ on that open cell. Both the candidate $v$ and the convolution have continuous representatives there, the latter by translation continuity in $L^2$. Vanishing at the $N$ nodes forces the residual to be zero on the entire cell. The union of these cells has full measure. Thus $v\in S_k$ and the projection gives the actual eigenfunction equation. No continuity across knots has to be assumed. If $v=0$ then $Q$ vanishes on the cells, so the whole unknown vector is zero. At $z=0$ the conditions would force $v=Q$ and $v\perp Q$, hence $v=0$, proving full column rank there. Real rank loss is equivalent to $\mathcal D_{k,a}(z)=0$. Compactness and the positive operator norm in \eqref{eq:odd-compression} prove the minimum exists and give its value. Reversing the isometry proves the assertion about extremizers.
\end{proof}

\section{One Fourier--Bessel sequence in every dimension}
\label{sec:ef-fourier-bessel}
Fix $d\ge2$, $0<a<1$, $R=1/2$, and choose $L=R+a$. Put $\nu=d/2-1$, $\Omega_d(t)=\Gamma(\nu+1)(2/t)^\nu J_\nu(t)$, with value one at zero. Let $\alpha_n$ be the positive zeros of $J_\nu$ in increasing order, $k_n=\alpha_n/L$, and, for $0<r<L$, set
\[
 e_n(r)=w_n r^{-\nu}J_\nu(k_n r),\qquad w_n=\frac{\sqrt2}{L|J_{\nu+1}(\alpha_n)|},\qquad
 m_n=\Omega_d(ak_n).
\]
These are the normalized radial Dirichlet eigenfunctions on the enclosing ball, in $L^2((0,L),r^{d-1}\,dr)$. Their completeness is the Fourier--Bessel expansion theorem \cite[Sec.~10.23(iii)]{DLMFBessel}. Write $u_n=e_n|_{(0,R)}$ and $G_N=(g_{ij})_{i,j\le N}$ for the real Gram matrix in $L^2((0,R),r^{d-1}dr)$. Its entries are explicit. Writing $t_i=k_iR$, $j_i=J_\nu(t_i)$, $h_i=J_{\nu+1}(t_i)$, the Lommel integrals give
\begin{align*}
 g_{ij}&=w_iw_j R\frac{k_i h_i j_j-k_j j_i h_j}{k_i^2-k_j^2}\quad(i\ne j),\\
 g_{ii}&=w_i^2\frac{R^2}{2}\left[j_i^2-J_{\nu-1}(t_i)h_i\right].
\end{align*}
The identities follow as well by integrating the Bessel differential equation for the two orders of frequency and passing to the coincident limit. Let $D_N=\operatorname{diag}(m_1,\ldots,m_N)$, $c_d=\Gamma(d/2)/(\sqrt\pi\,\Gamma((d-1)/2))$ and
\[
 \eta_N=\min\{1,8c_d(ak_{N+1})^{-1/2}\}.
\]
For $d\ge3$ the smaller bound $\min\{1,4c_d/(ak_{N+1})\}$ can also be used. These bounds follow by one-dimensional angular integration: $\Omega_d$ is the Fourier transform of $c_d(1-v^2)^{(d-3)/2}$ on $(-1,1)$. For $d\ge3$ integration by parts on the two monotone halves gives $4c_d/t$ at $t>0$. For $d=2$ remove endpoints of length $t^{-1}$ when $t\ge1$. Their combined mass is at most $4c_d/\sqrt{t}$ and bounded variation on the remainder contributes at most the same. This proves $|m_n|\le\eta_N$ for $n>N$, with the displayed $t^{-1/2}$ estimate valid for all $d\ge2$.

For a real diagonal $N$-by-$N$ matrix $D$, write $\Lambda_N(D)=\max\{0,\lambda_{\max}(G_N^{1/2}DG_N^{1/2})\}$. Then
\begin{equation}\label{eq:ef-bessel-enclosures}
 -\eta_N+\Lambda_N(D_N+\eta_N I)\ \le\ \Phi_d(a)\ \le\
 \eta_N+\Lambda_N(D_N-\eta_N I).
\end{equation}
In particular the interval width is at most $2\eta_N$. It often shrinks much faster. The following scalar entire function gives one equation for every dimension, including each even dimension,
\begin{equation}\label{eq:ef-bessel-secular}
 \mathcal D_{d,a}(z)=\lim_{N\to\infty}\det\!\left[I_N-z(D_NG_N)^4\right],
 \qquad \Phi_d(a)=\left(\min\{z>0:\mathcal D_{d,a}(z)=0\}\right)^{-1/4}.
\end{equation}
The limit is locally uniform in $z$. Each approximating determinant uses a single sequence of $N$ frequencies, independently of the ambient dimension. This is an infinite determinant in even dimensions and makes no claim of an elementary expression in terms of finitely many classical functions.

To prove the formulas, extend a radial half-ball function by zero to the enclosing ball, by the isometry $E$. The restriction to the half-ball of the spherical average of each $e_n$ is $m_n u_n$. Indeed, its entire radial Helmholtz extension is a constant times the spherical average of the plane waves $e^{ik_n x\cdot\theta}$. Averaging these waves over a sphere of radius $a$ multiplies them by $m_n$, and the integration uses only radii at most $R+a=L$. Finite Fourier--Bessel sums are dense, and the restricted spherical average is a contraction from the enclosing-ball radial space to the inner one by zero extension and Young's inequality. Therefore the radial spherical averaging operator $A$ is exactly the compression by $E$ of the diagonal multiplier $\operatorname{diag}(m_n)$ in the enclosing-ball basis. In particular
\[
 A_N=\sum_{n\le N}m_n u_n\langle\,\cdot\,,u_n\rangle,qquad \|A-A_N\|\le\eta_N.
\]
The upper and lower operator bounds replace every omitted diagonal entry by $\pm\eta_N$, respectively. Since $\sum_nu_n\langle\,\cdot\,,u_n\rangle=I$ strongly, and the nonzero eigenvalues of a finite sum with diagonal coefficients $D$ equal those of $G_N^{1/2}DG_N^{1/2}$, this gives \eqref{eq:ef-bessel-enclosures}. Positivity of the physical kernel implies $\Phi_d(a)=\sup\sigma(A)=\|A\|$.

For completeness the trace-ideal argument justifying the entire limit is as follows. The zeros satisfy $\alpha_n\ge n$. Indeed the defining alternating Bessel series is positive on $(0,1]$ after removal of its positive leading factor. For $x\ge1$, $\sqrt{x}J_\nu(x)$ solves $y''+(1-(\nu^2-1/4)/x^2)y=0$ with coefficient at most $5/4$. Integration on a nodal interval and the Dirichlet Poincar\'e inequality give a length of at least $2\pi/\sqrt5>1$. Thus $|m_n|\le C_{d,a} n^{-1/2}$ for $C_{d,a}=8c_d\sqrt{L/a}$. Compression by an isometry cannot increase Schatten norms. In particular $\|A-A_N\|_4\le C_{d,a}N^{-1/4}$ and $\max(\|A\|_4,\|A_N\|_4)\le C_{d,a}(\sum_{n=1}^\infty n^{-2})^{1/4}$. The fourth powers therefore converge in trace norm, by telescoping and the Schatten H\"older inequality. The trace-class Fredholm determinant is continuous in trace norm, locally uniformly on the $z$-plane \cite[Chap.~3]{SimonTrace}. For each $N$ the nonzero eigenvalues of $A_N$ with multiplicities are those of $D_NG_N$. Consequently \eqref{eq:ef-bessel-secular} equals $\det(I-zA^4)$. The nonzero positive compact operator $A^4$ has largest eigenvalue $\|A\|^4$, proving the stated first-root formula.

\section{The small-radius limit}
\label{sec:smallradius}
For each fixed $d\ge2$ there is a simple asymptotic expression even when the finite matrices grow as $a\downarrow0$. If $j_{\nu,1}$ is the first positive zero of the ordinary Bessel function $J_\nu$, then
\begin{equation}\label{eq:small-radius}
 \Phi_d(a)=1-\frac{2j_{d/2-1,1}^2}{d}\,a^2+o(a^2).
\end{equation}
In particular $(1-\Phi_3(a))/a^2\to2\pi^2/3$. The constant is related to the minimum-variance theorem of Ehm, Gneiting and Richards \cite[Theorem~5.1]{EGR}. Their inequality concerns the curvature of a fixed characteristic function at zero. To pass to the envelope one must control optimizing functions that vary with $a$.

Here $B$ is the half-ball of Section~\ref{sec:sered}, and $j_\nu(t)=\Gamma(\nu+1)(2/t)^\nu J_\nu(t)$ is the normalized Bessel function. Use the half-ball operator $A$ of that section. Its kernel is positivity preserving, so $|\langle Ah,h\rangle|\le\langle A|h|,|h|\rangle$ and its norm is the supremum of its positive Rayleigh quotients. With unitary Fourier transformation, Plancherel therefore gives
\[
 \frac{1-\Phi_d(a)}{a^2}=
 \inf_{\substack{\|h\|_2=1\\\operatorname{supp}h\subseteq B}}
 \int_{\mathbb R^d} Q_a(\xi)|\widehat h(\xi)|^2\,d\xi,
 \qquad Q_a(\xi)=\frac{1-j_{d/2-1}(a|\xi|)}{a^2}.
\]
Since $j_{d/2-1}(t)$ is the characteristic function of the first coordinate of a uniform sphere point, its variance is $1/d$. The inequality $1-\cos u\le u^2/2$ gives $0\le Q_a(\xi)\le |\xi|^2/(2d)$, and $Q_a\to |\xi|^2/(2d)$ uniformly on compact frequency sets. Testing on the normalized first Dirichlet eigenfunction on $B$ and using dominated convergence gives a limsup at most $\lambda_1(B)/(2d)$.

To prove the reverse inequality, the coordinate density is continuous in distribution, so $j_{d/2-1}(t)<1$ for $t>0$ and tends to zero at infinity. Its expansion at zero and compactness of intervening intervals imply $1-j_{d/2-1}(t)\ge c\min(t^2,1)$ for a constant $c>0$. Given any sequence $a_n\downarrow0$, take normalized trial functions within $1/n$ of the respective infima. Their energies are bounded by the preceding upper bound, so for every fixed $M>0$ and all large $n$ their squared Fourier tails outside $|\xi|\le M$ are bounded by $C/M^2$. Together with the fixed spatial support, these tail estimates give $L^2$ precompactness by the Fr\'echet--Kolmogorov criterion. Pass to a strongly convergent subsequence realizing the liminf, with limit $h$ of norm one supported in $B$. Restriction to a frequency ball, uniform convergence of $Q_{a_n}$ there and then monotone passage $M\to\infty$ give a liminf at least $\|\nabla h\|_2^2/(2d)$. In particular $h\in H^1(\mathbb R^d)$, with zero extension from the open ball, so $h|_{\operatorname{int} B}\in H_0^1(\operatorname{int} B)$. The Dirichlet variational inequality now gives the reverse bound. Since $B$ has radius $1/2$, its first eigenvalue is $4j_{d/2-1,1}^2$, proving \eqref{eq:small-radius}. The same compactness argument shows that normalized nonnegative radial maximizing factors converge in $L^2$ to the positive normalized first Dirichlet eigenfunction. Simplicity of that eigenvalue identifies every limit point.

\section{The sharp high-dimensional asymptotic}
\label{sec:ef-high-d}
Fix $0<a<1$. For this section $R=1/2$, $p=(d-3)/2$, $M=1-a^2$ and $\gamma=2a^2/M$. Then
\begin{equation}\label{eq:ef-high-d}
 \Phi_d(a)\sim \frac{(1-a^2)^{(d-1)/2}}{2a^3\sqrt{2\pi d}}
 \qquad(d\longrightarrow\infty).
\end{equation}
There is also an asymptotic shape for the radial factors: in the unweighted radial coordinate, on the scale $x=p(R-r)$, their normalized nonnegative profiles tend in $L^2(0,\infty)$ to $(2\gamma)^{1/2}e^{-\gamma x}$, after extending the rescaled profiles by zero past $pR$.

To prove the claims, restrict to $d>3$. Under the radial isometry $u(r)\mapsto r^{(d-1)/2}u(r)$, with a fixed surface-area normalization if needed, the spherical average restricted to the half-ball has symmetric kernel on $(0,R)$
\begin{equation}\label{eq:trianglekernel}
 K_d(r,s)=\frac{c_d}{a}\, X(r,s)^p,\qquad c_d=\frac{\Gamma(d/2)}{\sqrt\pi\,\Gamma((d-1)/2)},\qquad
 X(r,s)=\frac{[4r^2s^2-(r^2+s^2-a^2)^2]_+}{4a^2rs}.
\end{equation}
Indeed the first coordinate of a uniform point of the unit sphere has density $c_d(1-v^2)^p$. Substituting $s^2=r^2+a^2+2arv$ into the spherical mean gives the kernel. In particular the norm of $K_d$ is $\Phi_d(a)$, by the radial extremal reduction. Define $X=0$ if $rs=0$. This extension is continuous because $0\le X\le rs/a^2$.

The unique maximizer of $X$ on the closed square is $(R,R)$, where it equals $M$. Indeed, on the positive region, for $r\ge s>0$ set $t=s/r$ and $b=a/r$. Then
\[
 4b^2 X=-t^3+2(1+b^2)t-(1-b^2)^2/t.
\]
The derivative of this expression is $-3t^2+2(1+b^2)+(1-b^2)^2/t^2$. It decreases with $t$ and at $t=1$ equals $b^4>0$. Increasing $s$ to $r$ preserves the strict triangle inequalities whenever they held initially. On the positive part of the diagonal, $X(r,r)=1-a^2/(4r^2)$ strictly increases with $r$. Symmetry proves the assertion. At $(R,R)$, both first partial derivatives equal $a^2/(4R^3)=2a^2$. There exists $c>0$ with
\[
 X(r,s)/M\le e^{-c(2R-r-s)}
 \qquad(0\le r,s\le R).
\]
This follows from the positive partial derivatives in a neighborhood of the corner and, on the compact complement, the strict maximum just proved.

Under the unitary rescaling $v(r)\mapsto p^{-1/2}v(R-x/p)$, and after division by $c_dM^p/(ap)$, the kernel becomes
\[
 \mathbf1_{(0,pR)^2}(x,y)\left[\frac{X(R-x/p,R-y/p)}M\right]^p.
\]
It tends pointwise to $e^{-\gamma(x+y)}$ for finite $x,y>0$ and is bounded by $e^{-c(x+y)}$. Dominated convergence gives Hilbert--Schmidt convergence, hence operator-norm convergence. The limit operator has rank one, with eigenvalue $1/(2\gamma)$. This proves $\Phi_d(a)\sim c_d M^p/(2a p\gamma)$. Stirling's formula yields \eqref{eq:ef-high-d}. Norm convergence to the rank-one limit also gives the maximizing profile claim: the squared projection of any unit nonnegative top eigenvector onto the normalized exponential tends to one by its Rayleigh quotient, and the sign is fixed by nonnegativity.

\section*{Relation to earlier results and priority}
The dated original questions identify the targets independently of the auxiliary constructions. For the Gaussian problem the closest antecedent is Theorem~3.2 of \cite{JMR2009}. Its proof pairs the nonreal zeros and bounds the degree of a polynomial quotient for a compactly supported radial function. Proposition~3.4 there gives an extreme compact example with nonreal zeros. Gaussian convolution changes this behavior. The critical-order sector estimate, the reverse construction of perturbations by polynomial division, and the realization of all the polynomial directions in Theorem~\ref{thm:faces} give a necessary and sufficient condition and exact face dimensions.

For Question~2, general Mellin methods for scale mixtures predate this work. The explicit family-wide spectral multiplier and control of escaping shape and scale parameters are the additional ingredients. The quartic parameter arcs come from \cite{JMR2009}. The signed criterion combines Gaussian testing with the classical Stieltjes moment theorem. For compact covariances, the one-dimensional factorization and countable multivariate expansions are prior results \cite{EGR,GM}. For the spherical problem the largest-eigenvalue reduction is due to Efimov and is treated in the wider framework of \cite{GM}. The finite equations, the matrix enclosures and the limiting estimates concern that original extremal value.

The literature check, completed 20 November 2026, compared the original papers, the direct citing records indexed by OpenAlex for the extremal-ray paper, and keyword and author searches in publication and preprint indexes including DataCite. Searches included the tent--Gaussian example, extreme positive positive definite functions, and the named kernel families. No earlier answer to the two parts of Question~3 or identical heat-flow face classification was located in those records. The arguments are presented for independent checking. The search does not establish absolute priority or comparative historical importance.

\begin{thebibliography}{99}
\bibitem{JMR2009} P. Jaming, M. Matolcsi and Sz. Gy. R\'ev\'esz, On the extremal rays of the cone of positive, positive definite functions, \emph{J. Fourier Anal. Appl.} \textbf{15} (2009), 561--582. \href{https://doi.org/10.1007/s00041-008-9057-6}{Journal article}; \href{https://arxiv.org/abs/0801.0941}{arXiv:0801.0941} (2008).
\bibitem{Borisov} A. Borisov, Positive positive-definite functions and measures on locally compact abelian groups, 1999. \href{https://arxiv.org/abs/math/9906126}{arXiv:math/9906126}.
\bibitem{EF} A. Eremenko and A. Fryntov, Stability in the Marcinkiewicz theorem, \emph{Zh. Mat. Fiz. Anal. Geom.} \textbf{17} (2021), 463--467. \href{https://arxiv.org/abs/2106.14078}{arXiv:2106.14078}.
\bibitem{Akhiezer} N. I. Akhiezer, \emph{The Classical Moment Problem and Some Related Questions in Analysis}, Oliver \& Boyd, 1965.
\bibitem{EGR} W. Ehm, T. Gneiting and D. Richards, Convolution roots of radial positive definite functions with compact support, \emph{Trans. Amer. Math. Soc.} \textbf{356} (2004), 4655--4685. \href{https://doi.org/10.1090/S0002-9947-04-03502-0}{Journal article}.
\bibitem{GM} D. V. Gorbachev and A. D. Manov, Siegel--Sz\'asz extremal problem for a power weight, \emph{Sb. Math.} \textbf{217} (2026), 342--368. \href{https://www.mathnet.ru/eng/sm10268}{Article and references}.
\bibitem{EfimovSphere} A. V. Efimov, A version of the Turan problem for positive definite functions of several variables, \emph{Trudy Inst. Mat. i Mekh. UrO RAN} \textbf{17} (2011), no.~3, 136--154. English translation in \emph{Proc. Steklov Inst. Math.} \textbf{277} (2012), suppl.~1, S93--S112. \href{https://doi.org/10.1134/S0081543812050100}{Journal article}.
\bibitem{Freeman} J. Freeman and S. J. Miller, Determining optimal test functions for bounding the average rank in families of $L$-functions, \href{https://arxiv.org/abs/1507.03598}{arXiv:1507.03598} (2015).
\bibitem{DLMFBessel} F. W. J. Olver et al., eds., \emph{NIST Digital Library of Mathematical Functions}, Chapter~10. \href{https://dlmf.nist.gov/10}{Bessel functions}.
\bibitem{SimonTrace} B. Simon, \emph{Trace Ideals and Their Applications}, second ed., AMS, 2005.
\bibitem{Corpus} QEEDDD, \emph{Verifiable Mathematical Reasoning Corpus}, entries 027, 029, 103 and 126. \href{https://github.com/QEEDDD/Verifiable-Mathematical-Reasoning-Corpus/tree/c2ac34ea4622778187173f41bb1996deee8081b7}{Repository snapshot c2ac34e}.
\bibitem{RTNewman} B. Rodgers and T. Tao, The de Bruijn--Newman constant is non-negative, \emph{Forum Math. Pi} \textbf{8} (2020), e6. \href{https://doi.org/10.1017/fmp.2020.6}{Journal article}.
\bibitem{DobnerNewman} A. Dobner, A proof of Newman's conjecture for the extended Selberg class, \emph{Acta Arith.} \textbf{201} (2021), 29--62. \href{https://arxiv.org/abs/2005.05142}{Preprint}.
\bibitem{KKLNewman} H. Ki, Y.-O. Kim and J. Lee, On the de Bruijn--Newman constant, \emph{Adv. Math.} \textbf{222} (2009), 281--306. \href{https://doi.org/10.1016/j.aim.2009.04.003}{Journal article}.
\bibitem{BombieriWeil} E. Bombieri, Remarks on Weil's quadratic functional in the theory of prime numbers, I, \emph{Rend. Mat. Acc. Lincei}, ser.~9, \textbf{11} (2000), 183--233. \href{https://www.bdim.eu/item?fmt=pdf&id=RLIN_2000_9_11_3_183_0}{Full text}.
\bibitem{CsordasXi} G. Csordas, Fourier transforms of positive definite kernels and the Riemann $\xi$-function, \emph{Comput. Methods Funct. Theory} \textbf{15} (2015). \href{https://doi.org/10.1007/s40315-014-0105-8}{Journal article}.
\end{thebibliography}
\end{document}
